\documentclass[letterpaper]{article} % DO NOT CHANGE THIS
\usepackage[submission]{aaai2027}  % DO NOT CHANGE THIS
\usepackage[hyphens]{url}  % DO NOT CHANGE THIS
\usepackage{graphicx} % DO NOT CHANGE THIS
\urlstyle{rm} % DO NOT CHANGE THIS
\def\UrlFont{\rm}  % DO NOT CHANGE THIS
\usepackage{natbib}  % DO NOT CHANGE THIS AND DO NOT ADD ANY OPTIONS TO IT
\usepackage{caption} % DO NOT CHANGE THIS AND DO NOT ADD ANY OPTIONS TO IT
\frenchspacing  % DO NOT CHANGE THIS

\usepackage{amsmath}
\usepackage{amssymb}
\usepackage{amsthm}

% ── main.tex 와의 병합 규약 ───────────────────────────────────────────────
% 병합 절차: (1) 이 블록(\usepackage{amsmath} 부터 "병합 규약 끝" 까지)을
% main.tex 프리앰블 끝에 그대로 붙인다. (2) supplement 의 나머지 프리앰블은 버리고
% \appendix 부터 이 파일의 \bibliography 직전까지를 main.tex 의 \bibliography
% 뒤, \end{document} 직전에 넣는다. (3) main.tex 의 \bibliography 와
% \end{document} 만 한 번 유지한다.
% 블록을 통째로 버리면 안 된다 -- main.tex 는 proposition 만 정의하고
% lemma/corollary 가 없어서 부록이 컴파일되지 않는다.
% 블록은 멱등하다. main.tex 가 이미 정의한 것(\ms, \msb, \vb, ourshade, \ourrow,
% B, \lval, proposition)은 건너뛰고 없는 것만 새로 만들므로, 붙여도 중복 정의
% 오류가 나지 않고 렌더 결과도 달라지지 않는다(본문 Table 1--4 / 부록 S1--S13,
% 정리 번호 모두 단독 컴파일과 같음을 확인). 정의를 바꿀 일이 생기면 main.tex 가
% 제출본이라 고칠 수 없으므로 반드시 main.tex 쪽 정의에 맞춘다.
%
% 표 스타일 헬퍼. colortbl 을 직접 부르는 이유: aaai2027.sty 가 xcolor 를 먼저
% 읽을 경우 \usepackage[table]{xcolor} 가 option clash 를 낸다. colortbl 은
% array 를 함께 읽어 오므로 >{...} 열 수식자와 \newcolumntype 도 쓸 수 있다.
% (main.tex 는 \PassOptionsToPackage 로 같은 일을 하므로 합쳐도 충돌하지 않는다.)
\usepackage{xcolor}
\usepackage{colortbl}
\providecolor{ourshade}{RGB}{228,238,250}     % 우리 방법 강조용 옅은 파랑
% AAAI 는 표 글씨의 하한을 9pt 로 규정하고(CameraReady2027 "Tables") 글자 크기를
% 바꾸는 명령 자체를 금지하므로, 표준편차도 표 본문(\small = 9pt)과 같은 크기로
% 둔다. main.tex 의 \ms/\msb 와 동일한 정의.
\providecommand{\ms}[2]{\ensuremath{#1 \pm #2}}   % 평균 ± 표준편차
\providecommand{\msb}[2]{\ensuremath{\mathbf{#1} \pm #2}}% 열 최고값(굵게) ± 표준편차
\providecommand{\vb}[1]{\ensuremath{\mathbf{#1}}} % 표준편차 없는 열 최고값
\providecommand{\ourrow}{\rowcolor{ourshade}}     % 행 배경
\makeatletter
\@ifundefined{NC@find@B}{%
  \newcolumntype{B}{>{\columncolor{ourshade}}c}% 배경이 깔린 가운데 정렬 열
}{}
\makeatother

\providecommand{\lval}{\ell_{\mathrm{val}}}

% 한 표에 여러 설정을 담을 때 쓰는 블록 제목 행(\multicolumn{N}{l}{\emph{...}})의
% 위아래 여백. 제목이 \hline 과 첫 데이터 행 사이에 그대로 끼면 어느 행이 제목인지
% 한눈에 안 들어와서, 위(\hline 다음)와 아래(제목 행 끝)에 이만큼씩 띄운다.
% booktabs 의 \addlinespace 와 같은 방식이며, 값을 바꾸면 모든 표에 함께 적용된다.
% main.tex 에는 이런 블록 제목 행이 없어 합본해도 본문 표는 달라지지 않는다.
\providecommand{\ourpanelsep}{2pt}

% main.tex 는 proposition 만 정의하므로, 합본과 단독 컴파일 모두에서 없는 환경만
% 추가한다. 본문에는 proposition 인스턴스가 없어 두 경우의 부록 번호도 동일하다.
\theoremstyle{definition}
\makeatletter
\@ifundefined{proposition}{\newtheorem{proposition}{Proposition}}{}
\@ifundefined{lemma}{\newtheorem{lemma}{Lemma}}{}
\@ifundefined{corollary}{\newtheorem{corollary}{Corollary}}{}
\makeatother
% ── 병합 규약 끝 ─────────────────────────────────────────────────────────

\pdfinfo{
/TemplateVersion (2027.1)
}

\setcounter{secnumdepth}{2}

\title{Supplementary Material for\\Flirds: Efficient Client-Level Contribution Evaluation in Federated Learning via In-Run Data Shapley}
\author{
    Anonymous Submission
}
\affiliations{
}

\begin{document}

\maketitle
\appendix

% 부록 표는 S 접두를 붙여 본문 Table 1--4 와 번호가 겹치지 않게 한다(내부 \ref 는
% 자동으로 S1, S2, ... 로 렌더링되고, 합본에서는 본문 표를 가리키는 \ref 도 본문
% 번호로 렌더링된다). 프리앰블이 아니라 \appendix 뒤에 두는 이유: 두 파일을
% 합치면 프리앰블에 있을 때 본문 Table 1--4 까지 S 접두가 붙는다. 카운터 초기화도
% 같은 이유이며, 별도 컴파일에서는 이미 0 이라 아무 영향이 없다.
\setcounter{table}{0}
\renewcommand{\thetable}{S\arabic{table}}

% ===================================================================== A
\section{Proofs and Formal Statements}
\label{app:proofs}

% 귀속은 선언문 없이 헤더 인용으로만 표시한다. Lemma 1 은 알려진 사실의
% 재진술이라 인용을 그대로 달고(증명은 자기완결성을 위해 유지), closed form 은
% IRDS 의 유도를 라운드 게임으로 옮긴 것이라 after \citealp 로 달며 Theorem 이
% 아니라 Proposition 으로 둔다(레이블 thm:closed 는 이전 검토 기록과의 일치를
% 위해 유지).

\subsection{System Model and Assumptions}
\label{app:model}

We assume synchronous FedAvg. All participants in $P_r$ begin local training
in round $r$ from the same model $w^r$ and submit the client updates
$\delta_k^r := w_k^r - w^r$. The server aggregates them as
$w^{r+1} = w^r + \sum_{k \in P_r} p_k^r \delta_k^r$ with
$p_k^r = n_k / \sum_{j \in P_r} n_j$. The weights are held fixed
independently of the coalition $S$. The realized trajectory we fix is
$\{w^r, \{\delta_k^r\}_{k \in P_r}\}_{r=0}^{R-1}$, together with the local
dataset sizes $\{n_k\}$ that determine the weights. We write
$g_r = \nabla\lval(w^r)$, $H_r = \nabla^2\lval(w^r)$ (the true Hessian),
$\Delta_S^r = \sum_{k \in S \cap P_r} p_k^r \delta_k^r$, and
$\Delta W_r = \Delta_{P_r}^r$. Because $\phi$ attributes a decrease in
validation loss, a beneficial client receives $\phi_k > 0$.

% 적용 범위: 실험은 FedAvg 지만 부록 A 가 소비하는 성질은 (i)-(iii) 뿐이므로,
% 이를 만족하는 집계 규칙에는 p_k^r -> c_k^r 치환만으로 그대로 성립한다.
% 알고리즘 인스턴스는 열거하지 않는다(적용 여부 판단은 독자 몫). main
% Limitations 의 "coalition-independent additive aggregation" 절을 조건으로
% 명시화한 것이다.
The experiments in this paper evaluate FedAvg. The analysis, however, uses
only three properties of this system. (i)~Rounds are synchronous, so the
updates received in round $r$ are applied together at the common current
model $w^r$. (ii)~The server observes each participant's update
$\delta_k^r$ individually and holds the validation set that defines
$\lval$. (iii)~The realized model recursion is
$w^{r+1} = w^r + \sum_{k \in P_r} c_k^r\, \delta_k^r$ with coefficients
held fixed at their realized values, independently of any coalition being
evaluated. The coefficients may
depend on the round, on the participating set, or on the realized updates
themselves. What is excluded is dependence on the coalition $S$
(Appendix~\ref{app:closed}). For any aggregation rule with these three
properties, Appendices~\ref{app:game}--\ref{app:remainder} hold with
$c_k^r$ in place of $p_k^r$.

The required system assumption is not a stateless local optimizer but
stateless, additive server aggregation. When server state affects transitions
between rounds, as with server momentum or Adam, the grand-coalition
displacement no longer equals the realized update $w^{r+1} - w^r$, and the
telescoping argument below fails.

\subsection{The Round Game and Its Decomposition}
\label{app:game}

Define the round game and whole-run game as
$u_r(S) := \lval(w^r) - \lval(w^r + \Delta_S^r)$ and
$U^{\mathrm{in}} := \sum_r u_r$. The in-run Shapley is the Shapley value of
$U^{\mathrm{in}}$. Flirds computes exactly the Shapley value of
$\hat U^{\mathrm{in}}$, which is the sum of
$\hat u_r(S) := -\langle g_r, \Delta_S^r\rangle
- \tfrac12\langle \Delta_S^r, H_r \Delta_S^r\rangle$
(the first-order version is $\hat u_r^{(1)}$). The only difference between
the two games is the per-round Taylor remainder. We write $\hat\phi_k$ and
$\hat\phi_k^{(1)}$ for the Shapley values of $\hat U^{\mathrm{in}}$ and of
$\sum_r \hat u_r^{(1)}$, respectively.

The player associated with participating client $k$ is its weighted update
$p_k^r\delta_k^r$. The client update $\delta_k^r$ itself accumulates that
client's local optimization steps. The round game therefore differs from the
IRDS step game \citep{wang2025data}, and its larger expansion displacement
affects the remainder analyzed in Appendix~\ref{app:remainder}. Moreover, the
expansion point $w^r$ depends on the realized trajectory, so the Shapley-axiom
claims in Section~4.3 of the main paper apply only to the whole-run game on this
realized trajectory. The retraining-based Shapley \citep{ghorbani2019data} instead
retrains from scratch using only $S$, thereby changing both the weights and the
trajectory. Section~5.2 of the main paper evaluates the relationship between
these two games empirically.

\begin{lemma}[null-player reduction, \citealp{wang2023note}]
\label{lem:null}
If $v(\emptyset) = 0$ and $v(S) = v(S \cap T)$ for every $S$, then
$\phi_i(v) = 0$ for $i \notin T$ and $\phi_i(v) = \phi_i^T(v|_T)$ for
$i \in T$.
\end{lemma}

\begin{proof}
Every marginal contribution of $i \notin T$ is zero. The marginals of
$i \in T$ depend only on $S \cap T$. A uniform permutation of the full player
set induces a uniform relative order on $T$, so the resulting value equals
the permutation form of the Shapley value for the restricted game.
\end{proof}

\begin{proposition}[round decomposition and Shapley efficiency]
\label{prop:decomp}
If $S \supseteq P_r$, then $u_r(S) = \lval(w^r) - \lval(w^{r+1})$. In
particular, $U^{\mathrm{in}}([N]) = \lval(w^0) - \lval(w^R)$, and the values
of all clients sum to the validation-loss decrease over the complete run.
Furthermore, for every client $k$,
\[
\phi_k(U^{\mathrm{in}}) = \sum_{r:\,k \in P_r} \phi_k^{P_r}(u_r|_{P_r}).
\]
\end{proposition}

\begin{proof}
If $S \supseteq P_r$, then $\Delta_S^r = \Delta W_r$, so
$u_r(S) = \lval(w^r) - \lval(w^{r+1})$. Summing this identity over rounds
telescopes to $U^{\mathrm{in}}([N]) = \lval(w^0) - \lval(w^R)$. Because
$U^{\mathrm{in}}(\emptyset) = 0$, the Shapley efficiency axiom equates the
sum of all values with this decrease. Furthermore,
$u_r(S) = u_r(S \cap P_r)$, so Shapley linearity and Lemma~\ref{lem:null}
give the last display.
\end{proof}

This decomposition gives nonparticipating clients a value of zero in each
round. It rests only on $u_r(\emptyset) = 0$ and $u_r(S) = u_r(S \cap P_r)$.
The surrogate game $\hat u_r$ shares both properties, because
$\Delta_\emptyset^r = 0$ and $\Delta_S^r$ depends only on $S \cap P_r$, so the
same reduction gives
$\phi_k(\hat U^{\mathrm{in}}) = \sum_{r:\,k \in P_r} \phi_k^{P_r}(\hat u_r|_{P_r})$.
The decomposition itself does not require fixed weights. Fixed weights
are required for the coalition-independent additive structure and the closed
form below.

\subsection{Closed Form and Fixed Weights}
\label{app:closed}

\begin{proposition}[closed form, after \citealp{wang2025data}]
\label{thm:closed}
If $\lval$ is twice differentiable at $w^r$ and $H_r$ is symmetric, then for
$k \in P_r$
\[
\phi_k(\hat u_r)
= -p_k^r \langle g_r, \delta_k^r\rangle
- \tfrac12 p_k^r \langle \delta_k^r, H_r \Delta W_r\rangle,
\]
which is Eq.~(6) of the main paper.
\end{proposition}

\begin{proof}
Writing $q_{ij} := p_i^r p_j^r\langle \delta_i^r, H_r\, \delta_j^r\rangle$,
with the round index $r$ suppressed in $q$,
\[
\hat u_r(S)
= \sum_{k \in S}\Big(-p_k^r\langle g_r, \delta_k^r\rangle - \tfrac12 q_{kk}\Big)
- \sum_{\{i,j\} \subseteq S,\; i \ne j} q_{ij}.
\]
The first term is an additive game. By Shapley linearity the second term splits
into one game per pair. A pair game is nonzero only on coalitions containing
both of its players, so every other player is a null player and receives zero,
and the Shapley symmetry and Shapley efficiency axioms divide the interaction
equally between its two players. Hence,
$\phi_k(\hat u_r) = -p_k^r\langle g_r, \delta_k^r\rangle
- \tfrac12 \sum_{j \in P_r} q_{kj}
= -p_k^r\langle g_r, \delta_k^r\rangle - \tfrac12 p_k^r\langle \delta_k^r, H_r \Delta W_r\rangle$.
\end{proof}

In the final equality, every client shares the vector $H_r \Delta W_r$, so one
HVP per round suffices. The expression also holds when
$|P_r| = 1$, and $\sum_k \phi_k(\hat u_r) = \hat u_r(P_r)$ follows
immediately from the decomposition above.

Renormalizing the coefficients of the remaining clients within a partial
coalition gives
$\widetilde\Delta_S^r
= \sum_{k \in S} \frac{n_k}{\sum_{j \in S} n_j}\, \delta_k^r$,
which in general is not expressible as a sum of coalition-independent vectors
$b_k$. Even for two clients with equal $n$, the singletons force
$b_k = \delta_k^r$ while the pair gives
$\widetilde\Delta_{\{1,2\}}^r = \tfrac12(\delta_1^r + \delta_2^r) \ne b_1 + b_2$.
Fixed weights are therefore not required for every FL valuation, but they are
required for this additive structure and the one-HVP closed form. The
semantics also differ: ``had only these terms been reflected in the realized
round'' versus ``had only these clients participated.''

\noindent\textbf{Zero-update clients.} If $\delta_k^r = 0$, the estimator takes an
inner product with a zero tensor, so its value is algebraically exactly zero
regardless of numerical determinism. An exact zero obtained by evaluating
coalitions instead rests on the forward pass being reproducible to the last bit.
The zero-update conclusions in Section~5.3 of the main paper and
Appendix~\ref{app:intervention} rely on the estimator's algebraic zero, which
needs no such assumption.

\subsection{Taylor Remainder and Error Propagation}
\label{app:remainder}

Let $\mathcal U_r$ be an open convex set containing every coalition segment
$\{w^r + t\Delta_S^r : t \in [0,1],\, S \subseteq P_r\}$ and put
\[
M_{2,r} := \sup_{w \in \mathcal U_r}\|\nabla^2\lval(w)\|_{\mathrm{op}}, \quad
D_r := \max_{S \subseteq P_r}\|\Delta_S^r\|,
\]
\[
M_{3,r} := \sup_{w \in \mathcal U_r,\, \|x\| = 1}
|\nabla^3\lval(w)[x,x,x]|.
\]
The norms are the Euclidean norms of the relevant training coordinates.

\begin{proposition}[utility remainder]
\label{prop:remainder}
If $\lval \in C^3(\mathcal U_r)$ then
\[
\begin{aligned}
|u_r(S) - \hat u_r(S)| \le \tfrac{M_{3,r}}{6}\|\Delta_S^r\|^3,
\\
\qquad |u_r(S) - \hat u_r^{(1)}(S)| \le \tfrac{M_{2,r}}{2}\|\Delta_S^r\|^2.
\end{aligned}
\]
\end{proposition}

\begin{proof}
Apply the integral form of the Taylor remainder to
$h(t) = \lval(w^r + t\Delta_S^r)$ and bound
$|h''(t)| \le M_{2,r}\|\Delta_S^r\|^2$ and
$|h'''(t)| \le M_{3,r}\|\Delta_S^r\|^3$.
\end{proof}

\begin{lemma}[Shapley error propagation]
\label{lem:prop}
For any two games $v$ and $v'$ on the same player set,
$|\phi_i(v) - \phi_i(v')| \le 2\max_S |v(S) - v'(S)|$.
\end{lemma}

\begin{proof}
Each marginal contribution in the difference game is bounded by the
right-hand side, and a Shapley value is a convex combination of those
marginals.
\end{proof}

\begin{corollary}[player-level bound]
\label{cor:player}
Applying Proposition~\ref{prop:remainder} and Lemma~\ref{lem:prop} to the
round decomposition of Proposition~\ref{prop:decomp},
\[
\begin{aligned}
|\phi_k(U^{\mathrm{in}}) - \hat\phi_k|
\le \sum_{r:\,k \in P_r} \tfrac{M_{3,r}}{3} D_r^3,
\\
\qquad
|\phi_k(U^{\mathrm{in}}) - \hat\phi_k^{(1)}|
\le \sum_{r:\,k \in P_r} M_{2,r} D_r^2.
\end{aligned}
\]
\end{corollary}

This bound establishes the order of the remainder rather than a tight
numerical value: its constants range over the entire coalition set and are not
estimated here. The bound is conditional on the $C^3$ assumption and therefore
applies to the LLM track with smooth activations
(Appendix~\ref{app:lora}). For the ReLU-based CNN track, the assumption holds
only piecewise. The fidelity claim on that track therefore rests on the
empirical results in Section~5.2 of the main paper and the direct measurement
in Appendix~\ref{app:remainder-measured}. That measurement enumerates the same
coalitions in one setting from each track, including the CNN setting of
Table~1 of the main paper, and reports the utility remainder in
Proposition~\ref{prop:remainder} directly.

\subsection{LoRA Coordinates}
\label{app:lora}

The game, estimator, and in-run Shapley for the LLM track are defined and
computed in the same LoRA-factor coordinates $\theta$. Write $w(\theta)$ for
the full network weights, in which each adapted layer takes the value
$W_0 + BA$, with the pretrained matrix $W_0$ frozen and the factors $B$ and
$A$ collected in $\theta$. Appendices~\ref{app:game}--\ref{app:remainder}
then hold with $\theta$ in place of $w$ and $\lval(w(\theta))$ in place of
$\lval(w)$. The composition of smooth activations with the polynomial map
$\theta \mapsto W_0 + BA$ is $C^3$, so the assumption of
Appendix~\ref{app:remainder} is preserved as well.

The second-order surrogate is coordinate-dependent. Because $BA$ is bilinear,
the additive structure of factor updates generally cannot be reinterpreted as
per-client additive displacements in full-weight space. The curvature of the
LoRA map is itself part of the Hessian. Consequently, the closed-form and
remainder claims in this paper are limited to the LoRA-factor coordinates.

% ===================================================================== B
\section{Experimental Protocol}
\label{app:protocol}

This appendix completes the setup of Section~5.1 of the main paper. That
section and the result-table captions already specify each setting's model,
dataset, partition, client count, participation rate, and round count.
Section~5.1 of the main paper also gives the corruption fraction shared across the main
settings. Four additional settings are detailed below. The
first two are the five-domain corpus at $N{=}5$ with one domain per client and
its Alpaca IID counterpart, both with full participation and $R{=}10$. The
third is an Alpaca setting at $N{=}5$ with full participation and $R{=}30$, run
at all three model scales. The fourth distributes the same five domains over
$N{=}100$ clients, whose runtimes are reported in Table~4 of the main paper. We
also provide the training and data configuration needed to rerun all settings
and reproduce the table entries.
Appendix~\ref{app:train} describes training,
Appendix~\ref{app:conditions} details the conditions, and
Appendix~\ref{app:data} specifies the data distribution and evaluation sets.

\subsection{Training Configuration}
\label{app:train}

Both tracks use plain SGD without momentum and a constant learning rate:
$10^{-3}$ for the LLM track and $0.01$ for the CNN track. Per round, the LLM
track performs 10 local optimization steps, whereas the CNN track performs 5
local epochs with a batch size of 64. The LLM local batch size decreases with
model size to fit memory: 16 at 1B, 8 at 3B, and 4 at 7B. The 1B and 3B models
are the instruction-tuned Llama-3.2-1B-Instruct and
Llama-3.2-3B-Instruct \citep{grattafiori2024llama}. The 7B model is
Llama-2-7b \citep{touvron2023llama}, a base model rather than an
instruction-tuned one. It is the 7B model standard in the federated LLM literature
\citep{ye2024openfedllm}. The LoRA \citep{hu2022lora} factors that the LLM
clients train and exchange have rank 16 and scaling factor 32.
The maximum sequence length is 768 in the two five-domain settings and in the
Alpaca setting at $N{=}5$ ($R{=}10$), and 512 in the main LLM setting and the
other two Alpaca settings.
For CIFAR-10 \citep{krizhevsky2009learning}, we use FedSVCNN from the FedSV
line of work \citep{wang2020principled}: two
$5{\times}5$ convolutional layers with 32 and 64 channels, followed by a
512-unit dense layer. For MNIST, we use LeNet5 \citep{lecun1998gradient}. All
settings are trained and evaluated in fp32 with PyTorch.

\noindent\textbf{Computing infrastructure.} We ran the experiments on two Linux
machines: a node with NVIDIA B200 GPUs (AMD EPYC 9365, 2.2\,TiB of memory,
Ubuntu 24.04.3 LTS) and a node with NVIDIA RTX 3090 GPUs (Intel Xeon Gold
5320, 64\,GiB of memory, Rocky Linux release 9.4). Together, the experiments
used roughly 2{,}100 GPU-hours: about 1{,}180 on the RTX 3090 node and about
900 on the B200 node. The software stack comprised Python 3.11, PyTorch 2.11
with CUDA 13.0, transformers 5.5, PEFT 0.19, TRL 1.2, and datasets 4.8.

\noindent\textbf{Server-side valuation.} We evaluate the validation gradient $g_r$ and
the Hessian-vector product $H_r \Delta W_r$ in Algorithm~1 of the main paper on
the entire server validation set specified in Appendix~\ref{app:data}, not on a
minibatch. The round game and its estimator therefore use the same validation
loss. One JVP applied to the validation-gradient function returns both
quantities in fp32. We process the validation set in chunks that fit device
memory. Chunking changes only how the sum is accumulated. Methods that
evaluate coalition utilities directly use the same validation set, so all
methods in a given setting, condition, and seed share one utility.

\noindent\textbf{Baseline hyperparameters.} The remaining baseline settings are as
follows. GTG-Shapley \citep{liu2022gtg} draws permutations within a round
until the running mean over the last 10 of them changes by less than $5\%$,
evaluating at least $\max(30, |P_r|)$ and at most
$\lceil 0.8 \cdot 2^{|P_r|} \rceil$ permutations, and FedSV uses
$\max(30, 2|P_r|)$ permutations per round. For both methods we set the two
utility-difference truncation thresholds (within a permutation and between
rounds) to zero, so a coalition evaluation is skipped only by the convergence
criterion above, and the truncation constant $c$ of Table~4 of the main paper
is set by that criterion rather than by a utility threshold. ComFedSV draws
$M = \max(10, \lceil N \ln N \rceil)$ permutations and completes the partially
populated utility matrix by rank-5 alternating least squares (60 iterations,
ridge coefficient $0.1$) \citep{fan2022improving}. It is defined for partial
participation: under full participation every visited coalition is evaluated
directly and no entry requires completion, so we score it in the
partial-participation settings but not in the full-participation ones.
ShapleyFL \citep{sun2023shapleyfl} smooths its per-round contributions with an
exponential moving average of coefficient $\beta{=}0.3$, and FedIF
\citep{tang2025lightweight} smooths its round influences with an exponential
moving average of rate $\gamma{=}0.3$, the value its authors use for CIFAR-10.
We use both values in every setting.

\noindent\textbf{Hyperparameter selection.} The values above are chosen once per track
and used unchanged for every method, condition, model scale, and seed. No
hyperparameter is tuned per method or per condition, and no search over
learning rates, LoRA ranks, or round counts selects the numbers reported in
Section~5 of the main paper. Each baseline keeps the hyperparameter values
published with it. Every method is therefore scored on the trajectory that
this single configuration produces, so a method's result reflects its
valuation rule rather than a separately tuned training setup.

\noindent\textbf{Seeds and determinism.} Each of the three seeds
$s \in \{0, 1, 2\}$ in Section~5.1 of the main paper is a single integer that
seeds Python, NumPy, and PyTorch at the start of a run. It determines the
training-data shuffle and partition, the participants sampled in each round,
the corrupted clients (Appendix~\ref{app:conditions}), the random streams
used by each corruption condition, and the Monte Carlo permutations used by
GTG-Shapley, FedSV, and ComFedSV. On the LLM track, the LoRA factors use a
common initialization seed before the run seed is applied, making their
initialization identical across the three seeds.

\subsection{Conditions and Corrupted Clients}
\label{app:conditions}

\textbf{clean} is the uncorrupted control for the four corruption
conditions below.

\noindent\textbf{answer-swap} (LLM) permutes the responses (including the derivation
and final answer) among a prescribed fraction of questions held by the same
client. The format remains valid, but the question--response correspondence
is broken, producing a realistic form of mislabeling. The fraction is $0.7$
in the main LLM setting and $1.0$ in the five-domain non-IID setting.

\noindent\textbf{zero-update free-rider} submits $\delta_k^r = 0$ without local
training.

\noindent\textbf{gradient noise} (CNN) adds i.i.d.\ Gaussian noise
$\mathcal{N}(0, 0.1^2)$ to every entry of the submitted update after normal
local training.

\noindent\textbf{label-flip} (CNN) replaces a fraction $0.7$ of a corrupted client's
labels with a label drawn uniformly from the other nine classes.

\noindent\textbf{Strengths and corrupted clients.} The corruption strength is a constant
in every setting and is never drawn from a distribution. Its value is stated
with each condition above. The corrupted clients, by contrast, are realized
differently across settings, and we report the realized count rather than the
nominal fraction wherever the two can differ. In the main CNN setting
($N{=}100$), label-flip designates each client independently as corrupted with
probability $0.4$, giving 39, 48, and 47 corrupted clients under
seeds 0 through 2. The two update-level conditions instead draw exactly
$40$ clients uniformly without replacement. The full-participation CNN
settings ($N{=}10$) draw
exactly 4 clients under every corruption condition, so all three of their
corruption conditions use the same corrupted clients. On the LLM track, the same index range is
corrupted across seeds:
clients $0$--$19$ of 50 in the main LLM setting ($40\%$), clients
$\{10, 30, 50, 70, 90\}$ of 100 in the $N{=}100$ five-domain setting ($5\%$),
and one of five clients in the five-domain non-IID setting ($20\%$). The Alpaca
settings are clean throughout. The run seed shuffles the GSM8K training set
before it is partitioned and draws the per-client domain mixture in the
$N{=}100$ five-domain setting, so the same index range holds different data in
each seed.

\subsection{Data Distribution and Evaluation Sets}
\label{app:data}

For GSM8K \citep{cobbe2021training}, we distribute the 7{,}473 official
training questions uniformly and IID across clients, assigning 149 questions
to each client and leaving the remainder unused. Training follows the official
split. The server validation set contains 200 of the 1{,}319 official test
questions, and exact match is measured on the remaining 1{,}119 questions. We
use greedy decoding without sampling, generate at most 256 new tokens, and
truncate prompts to 512 tokens. A second LLM corpus pools five domains: medical
flashcards (Medical Meadow) \citep{han2023medalpaca}, legal QA
\citep{legalqa}, finance QA
(FiQA) \citep{maia201818}, math word problems with rationales (AQUA-RAT)
\citep{ling2017program}, and general instruction following (Dolly)
\citep{conover2023free}. We convert all five
public corpora to a common free-form
instruction--response format so that the shared validation loss is comparable
across domains. Each domain contributes 200 training, 20 validation, and 40 test
examples, and the validation and test portions are pooled on the server, which
gives a 100-example server validation set and a 200-example test set. At
$N{=}5$ each client holds the training examples of one domain. A second setting,
whose valuation cost is reported in Table~4 of the main paper, draws on the same five
domains but spreads them over
$N{=}100$ clients with a Dirichlet($\alpha{=}0.5$) mixture over domains per
client, so each client holds a skewed blend rather than a single domain.

Three settings partition Alpaca-GPT4 \citep{peng2023instruction} IID. The first
matches the size of the five-domain non-IID setting so that only the corpus and
partition differ: $N{=}5$, full participation, $R{=}10$, 200 training examples
per client, a 100-example server validation set, and 200 test examples. The
other two use Alpaca-GPT4 20k, hold out 200 validation and 1{,}000 test
examples, and divide the remainder among clients. They differ only in client
count, participation rate, and round count. The $N{=}5$ setting uses full
participation for $R{=}30$. The model-scale axis in Figure~2 of the main paper
uses $N{=}20$, 2/20 participation, and $R{=}200$.

Of the two CNN partitions of Section~5.1 of the main paper, the IID partition
is a uniform random partition, whereas the Dirichlet($\alpha{=}1$) partition
\citep{hsu2019measuring} skews label distributions and client dataset sizes
jointly. Either partition divides the official training split among the
clients, and the run seed determines it. Both datasets carry a
10{,}000-example official test split, which we divide into a 2{,}000-example
server validation set and 8{,}000 test examples. We report held-out test
accuracy for the CNN track.


\begin{table*}[!t]
\centering
\small
\setlength{\tabcolsep}{4pt}
% mean 행은 조건 행들의 집계이므로 \cline 으로 분리하고(두 패널 사이의 \qquad 는
% 비워 둔다), 굵게는 mean 행에만 쓴다.
% 오른쪽 패널은 네 블록이 모두 Spearman 이므로 지표를 패널 헤더에서 한 번만 밝히고
% 블록 헤더에는 partition 만 남긴다. 왼쪽은 두 지표가 번갈아 나오므로 블록 헤더가
% 지표를 계속 진다 -- 이렇게 해야 같은 줄에 "MNIST / Dir(1)" 이 두 번 나오면서
% 지표만 다른 상태가 사라진다.
\begin{tabular}{lcB@{\qquad}lcB}
\hline
\multicolumn{3}{l}{\emph{$N{=}100$, 10/100, $R{=}120$}} & \multicolumn{3}{l}{\emph{$N{=}10$, full participation, $R{=}10$; Spearman's $\rho$}} \\
Condition & Flirds-1st & Flirds & Condition & Flirds-1st & Flirds \\
\hline
\noalign{\vskip\ourpanelsep}
\multicolumn{3}{l}{\emph{Spearman's $\rho$, MNIST / Dir(1)}} & \multicolumn{3}{l}{\emph{CIFAR-10 / Dir(1)}} \\[\ourpanelsep]
clean & \ms{0.809}{0.065} & \ms{0.981}{0.017} & clean & \ms{0.931}{0.037} & \ms{0.992}{0.007} \\
zero-update & \ms{0.913}{0.048} & \ms{0.994}{0.005} & zero-update & \ms{0.987}{0.022} & \ms{1.000}{0.000} \\
gradient noise & \ms{0.264}{0.080} & \ms{0.939}{0.030} & gradient noise & \ms{-0.188}{0.280} & \ms{0.915}{0.064} \\
label-flip & \ms{0.990}{0.003} & \ms{0.989}{0.010} & label-flip & \ms{0.980}{0.025} & \ms{0.996}{0.007} \\
\cline{1-3}\cline{4-6}
mean & $0.744$ & \vb{0.976} & mean & $0.678$ & \vb{0.976} \\
\hline
\noalign{\vskip\ourpanelsep}
\multicolumn{3}{l}{\emph{Spearman's $\rho$, MNIST / IID}} & \multicolumn{3}{l}{\emph{CIFAR-10 / IID}} \\[\ourpanelsep]
clean & \ms{0.952}{0.020} & \ms{0.998}{0.001} & clean & \ms{0.535}{0.196} & \ms{0.952}{0.012} \\
zero-update & \ms{0.966}{0.006} & \ms{0.996}{0.006} & zero-update & \ms{0.901}{0.049} & \ms{0.978}{0.015} \\
gradient noise & \ms{0.362}{0.025} & \ms{0.906}{0.027} & gradient noise & \ms{0.139}{0.116} & \ms{0.968}{0.019} \\
label-flip & \ms{0.981}{0.003} & \ms{0.975}{0.023} & label-flip & \ms{0.952}{0.032} & \ms{0.996}{0.007} \\
\cline{1-3}\cline{4-6}
mean & $0.815$ & \vb{0.969} & mean & $0.632$ & \vb{0.974} \\
\hline
\noalign{\vskip\ourpanelsep}
\multicolumn{3}{l}{\emph{Pearson's $r$, MNIST / Dir(1)}} & \multicolumn{3}{l}{\emph{MNIST / Dir(1)}} \\[\ourpanelsep]
clean & \ms{0.692}{0.270} & \ms{0.965}{0.052} & clean & \ms{0.947}{0.039} & \ms{0.976}{0.012} \\
zero-update & \ms{0.843}{0.191} & \ms{0.983}{0.028} & zero-update & \ms{0.996}{0.007} & \ms{1.000}{0.000} \\
gradient noise & \ms{-0.027}{0.229} & \ms{0.950}{0.014} & gradient noise & \ms{0.269}{0.355} & \ms{0.972}{0.019} \\
label-flip & \ms{0.987}{0.006} & \ms{0.994}{0.005} & label-flip & \ms{1.000}{0.000} & \ms{1.000}{0.000} \\
\cline{1-3}\cline{4-6}
mean & $0.624$ & \vb{0.973} & mean & $0.803$ & \vb{0.987} \\
\hline
\noalign{\vskip\ourpanelsep}
\multicolumn{3}{l}{\emph{Pearson's $r$, MNIST / IID}} & \multicolumn{3}{l}{\emph{MNIST / IID}} \\[\ourpanelsep]
clean & \ms{0.966}{0.008} & \ms{0.998}{0.000} & clean & \ms{0.778}{0.110} & \ms{0.814}{0.165} \\
zero-update & \ms{0.986}{0.002} & \ms{1.000}{0.000} & zero-update & \ms{0.931}{0.061} & \ms{0.966}{0.049} \\
gradient noise & \ms{0.162}{0.085} & \ms{0.949}{0.012} & gradient noise & \ms{0.430}{0.352} & \ms{0.903}{0.067} \\
label-flip & \ms{0.991}{0.002} & \ms{0.995}{0.002} & label-flip & \ms{0.976}{0.021} & \ms{1.000}{0.000} \\
\cline{1-3}\cline{4-6}
mean & $0.776$ & \vb{0.986} & mean & $0.779$ & \vb{0.921} \\
\hline
\end{tabular}
\caption{Agreement with the in-run Shapley on the CNN track.
\textbf{Left:} MNIST, both partitions ($N{=}100$, 10/100 participation,
$R{=}120$). \textbf{Right:} CIFAR-10 and MNIST, both partitions ($N{=}10$,
full participation, $R{=}10$, 4/10 corrupted).}
\label{tab:inrun-fidelity}
\end{table*}


\begin{table*}[!t]
\centering
\small
\setlength{\tabcolsep}{4pt}
% 두 번째 헤더 행이 "Flirds-1st & Flirds" 를 그대로 두 번 반복하므로, 어느 쌍이 어느
% 지표에 속하는지 \cline 으로 묶어 준다(tab:inrun-fidelity / tab:target-stability
% / tab:gate-pr 과 같은 규약).
\begin{tabular}{lcBcB}
\hline
 & \multicolumn{2}{c}{relative $L_2$ error} & \multicolumn{2}{c}{sign agreement (\%)} \\
\cline{2-3}\cline{4-5}
Condition & Flirds-1st & Flirds & Flirds-1st & Flirds \\
\hline
\noalign{\vskip\ourpanelsep}
\multicolumn{5}{l}{\emph{CIFAR-10 / Dir(1)}} \\[\ourpanelsep]
clean & \ms{0.199}{0.164} & \msb{0.067}{0.046} & \msb{96.7}{5.8} & \msb{96.7}{5.8} \\
zero-update & \ms{0.087}{0.017} & \msb{0.014}{0.003} & \msb{100.0}{0.0} & \msb{100.0}{0.0} \\
gradient noise & \ms{2.081}{0.056} & \msb{0.541}{0.100} & \ms{60.0}{0.0} & \msb{100.0}{0.0} \\
label-flip & \ms{0.074}{0.014} & \msb{0.013}{0.007} & \msb{100.0}{0.0} & \msb{100.0}{0.0} \\
\hline
\noalign{\vskip\ourpanelsep}
\multicolumn{5}{l}{\emph{CIFAR-10 / IID}} \\[\ourpanelsep]
clean & \ms{0.154}{0.047} & \msb{0.131}{0.005} & \msb{100.0}{0.0} & \msb{100.0}{0.0} \\
zero-update & \ms{0.053}{0.020} & \msb{0.011}{0.006} & \msb{100.0}{0.0} & \msb{100.0}{0.0} \\
gradient noise & \ms{1.826}{0.140} & \msb{0.408}{0.042} & \ms{60.0}{0.0} & \msb{100.0}{0.0} \\
label-flip & \ms{0.070}{0.008} & \msb{0.017}{0.003} & \msb{100.0}{0.0} & \msb{100.0}{0.0} \\
\hline
\noalign{\vskip\ourpanelsep}
\multicolumn{5}{l}{\emph{MNIST / Dir(1)}} \\[\ourpanelsep]
clean & \ms{0.430}{0.132} & \msb{0.298}{0.160} & \msb{100.0}{0.0} & \msb{100.0}{0.0} \\
zero-update & \ms{0.223}{0.037} & \msb{0.066}{0.057} & \msb{100.0}{0.0} & \msb{100.0}{0.0} \\
gradient noise & \ms{1.345}{0.674} & \msb{0.618}{0.315} & \ms{73.3}{15.3} & \msb{86.7}{15.3} \\
label-flip & \ms{0.168}{0.050} & \msb{0.085}{0.045} & \msb{100.0}{0.0} & \msb{100.0}{0.0} \\
\hline
\noalign{\vskip\ourpanelsep}
\multicolumn{5}{l}{\emph{MNIST / IID}} \\[\ourpanelsep]
clean & \ms{0.654}{0.057} & \msb{0.394}{0.032} & \msb{100.0}{0.0} & \msb{100.0}{0.0} \\
zero-update & \ms{0.275}{0.031} & \msb{0.109}{0.056} & \msb{100.0}{0.0} & \msb{100.0}{0.0} \\
gradient noise & \ms{1.009}{0.231} & \msb{0.730}{0.106} & \ms{66.7}{5.8} & \msb{93.3}{5.8} \\
label-flip & \ms{0.194}{0.017} & \msb{0.104}{0.029} & \msb{100.0}{0.0} & \msb{100.0}{0.0} \\
\hline
\end{tabular}
\caption{Value-level agreement with the in-run Shapley (CIFAR-10 and MNIST,
both partitions; $N{=}10$, full participation, $R{=}10$, 4/10 corrupted).}
\label{tab:value-level}
\end{table*}

\begin{table*}[!t]
\centering
\small
\setlength{\tabcolsep}{2pt}
% mean 행은 조건 행들의 집계이므로 가로선으로 분리하고, 굵게는 mean 행에만 쓴다.
\begin{tabular}{lccccccB}
\hline
Condition & in-run SV & GTG-Shapley & FedSV & ShapleyFL & FedIF & Flirds-1st & Flirds \\
\hline
\noalign{\vskip\ourpanelsep}
\multicolumn{8}{l}{\emph{CIFAR-10 / IID}} \\[\ourpanelsep]
clean & \ms{-0.273}{0.252} & \ms{-0.196}{0.217} & \ms{-0.176}{0.503} & \ms{0.063}{0.476} & \ms{0.067}{0.164} & \ms{-0.131}{0.316} & \ms{-0.232}{0.220} \\
zero-update & \ms{0.625}{0.145} & \ms{0.638}{0.084} & \ms{0.286}{0.197} & \ms{-0.562}{0.122} & \ms{0.713}{0.228} & \ms{0.759}{0.043} & \ms{0.579}{0.159} \\
gradient noise & \ms{0.653}{0.118} & \ms{0.705}{0.039} & \ms{0.661}{0.063} & \ms{0.657}{0.050} & \ms{0.745}{0.092} & \ms{-0.204}{0.305} & \ms{0.644}{0.133} \\
label-flip & \ms{0.705}{0.031} & \ms{0.709}{0.061} & \ms{0.701}{0.079} & \ms{0.717}{0.102} & \ms{0.701}{0.035} & \ms{0.689}{0.091} & \ms{0.709}{0.032} \\
\cline{1-8}
mean & $0.428$ & $0.464$ & $0.368$ & $0.219$ & \vb{0.557} & $0.278$ & $0.425$ \\
\hline
\noalign{\vskip\ourpanelsep}
\multicolumn{8}{l}{\emph{MNIST / Dir(1)}} \\[\ourpanelsep]
clean & \ms{0.632}{0.334} & \ms{0.838}{0.060} & \ms{0.798}{0.110} & \ms{0.491}{0.176} & \ms{0.285}{0.242} & \ms{0.620}{0.396} & \ms{0.600}{0.378} \\
zero-update & \ms{0.915}{0.072} & \ms{0.968}{0.025} & \ms{0.952}{0.021} & \ms{-0.115}{0.431} & \ms{0.861}{0.051} & \ms{0.907}{0.066} & \ms{0.915}{0.072} \\
gradient noise & \ms{0.903}{0.036} & \ms{0.851}{0.077} & \ms{0.903}{0.117} & \ms{0.867}{0.061} & \ms{0.681}{0.062} & \ms{0.301}{0.353} & \ms{0.879}{0.024} \\
label-flip & \ms{0.960}{0.070} & \ms{0.960}{0.039} & \ms{0.931}{0.043} & \ms{0.851}{0.138} & \ms{0.564}{0.138} & \ms{0.960}{0.070} & \ms{0.960}{0.070} \\
\cline{1-8}
mean & $0.853$ & \vb{0.904} & $0.896$ & $0.524$ & $0.598$ & $0.697$ & $0.839$ \\
\hline
\noalign{\vskip\ourpanelsep}
\multicolumn{8}{l}{\emph{MNIST / IID}} \\[\ourpanelsep]
clean & \ms{0.438}{0.102} & \ms{0.192}{0.286} & \ms{-0.111}{0.452} & \ms{0.624}{0.231} & \ms{0.741}{0.150} & \ms{0.519}{0.335} & \ms{0.358}{0.099} \\
zero-update & \ms{0.717}{0.140} & \ms{0.754}{0.068} & \ms{0.721}{0.145} & \ms{0.452}{0.198} & \ms{0.813}{0.157} & \ms{0.675}{0.101} & \ms{0.667}{0.095} \\
gradient noise & \ms{0.754}{0.091} & \ms{0.729}{0.031} & \ms{0.669}{0.043} & \ms{0.810}{0.115} & \ms{0.770}{0.063} & \ms{0.337}{0.398} & \ms{0.745}{0.044} \\
label-flip & \ms{0.762}{0.086} & \ms{0.669}{0.061} & \ms{0.697}{0.087} & \ms{0.867}{0.053} & \ms{0.842}{0.021} & \ms{0.745}{0.087} & \ms{0.762}{0.086} \\
\cline{1-8}
mean & $0.668$ & $0.586$ & $0.494$ & $0.688$ & \vb{0.792} & $0.569$ & $0.633$ \\
\hline
\end{tabular}
\caption{Spearman's $\rho$ against the retraining-based Shapley (CIFAR-10 IID
and both MNIST partitions; $N{=}10$, full participation, $R{=}10$, 4/10
corrupted).}
\label{tab:grid-retrain}
\end{table*}

% ===================================================================== C
\section{Full Fidelity Results}
\label{app:fidelity}

\textbf{Table conventions.} As in Section~5.1 of the main paper, tables with a per-seed
spread report the mean $\pm$ standard deviation over the three seeds in
Section~5.1 of the main paper. Tables that instead aggregate over seed pairs, coalitions, or
table rows identify that unit in the caption. Every table abbreviates Shapley
value as SV.

Bold marks the best value along each comparison axis: the mean row or column
when present, each condition otherwise, or each reported quantity when the
table has neither. All ties are bold. The shaded row or column is Flirds. A
mean row or column is the unweighted average of the displayed condition means
and has no standard deviation. An em dash marks a quantity undefined for that
row or condition. The reference entries are the ``in-run SV'' row or column and
the vanilla, oracle-exclusion, and selection-random rows. They are not
comparison targets and carry no bold. We likewise do not rank a quantity that a method
can lower simply by acting less often.
ComFedSV appears only in the partial-participation settings, for the reason
given in Appendix~\ref{app:train}.

Section~5.2 of the main paper reports both CIFAR-10 partitions in the main CNN
setting and, under full participation, the CIFAR-10 Dirichlet partition.
This appendix reports the corresponding MNIST results, the
remaining partitions, and the LLM settings. Fidelity is measured throughout with Spearman's $\rho$
for rankings and Pearson's $r$ for values, the same two metrics as in
Section~5.2 of the main paper. Both are scale-free, so
Appendix~\ref{app:value-level} adds scale-sensitive value-level metrics for the
one comparison in which the estimate and its reference share units.
Only Flirds and Flirds-1st target the game in Eq.~(5) of the main paper. Their
correlations with the in-run Shapley can therefore be interpreted as
estimation fidelity, and for that reason Table~\ref{tab:inrun-fidelity}
scores the two variants alone against it, as Figure~2 of the main paper does.
The comparison that includes the other four methods therefore uses the
retraining-based Shapley (Tables~\ref{tab:grid-retrain}
and~\ref{tab:llm-retrain}), the common external
reference that no evaluated method targets directly. A low correlation there
does not by itself indicate a failure of those methods.


\subsection{Agreement with the In-Run Shapley}
\label{app:cnn-main}

Table~\ref{tab:inrun-fidelity} collects both CNN comparisons that use the
in-run Shapley as the reference. The left panel is the main CNN setting
($N{=}100$, 10/100 participation, $R{=}120$). Because
Figure~2 of the main paper already reports CIFAR-10 there, we give Spearman's
$\rho$ and Pearson's $r$ for the two MNIST partitions only. The right panel covers full
participation ($N{=}10$, $R{=}10$), where the in-run Shapley is enumerated
exhaustively over $2^{10}$ coalitions, and reports Spearman's $\rho$ for all
four dataset--partition combinations.

On MNIST in the main CNN setting, Flirds attains high rank and value agreement
in all four conditions of both partitions. Gradient noise is the hardest
condition in both. There, Flirds-1st drops sharply in rank agreement, and at the
level of values its correlation even reverses sign under the Dirichlet
partition. Averaged over the four conditions, Flirds leads it by a wide margin
under both partitions. The conclusion about the
second-order curvature term from Section~5.2 of the main paper therefore persists on MNIST.

The right panel repeats the pattern under full participation. Flirds leads
Flirds-1st in 15 of the 16 rows and ties it in the sixteenth (label-flip
under the MNIST Dirichlet partition), and under gradient noise Flirds-1st
alone reverses sign under the CIFAR-10 Dirichlet partition. The conclusion
from the partial-participation setting therefore also holds under full
participation.

\subsection{Value-Level Agreement}
\label{app:value-level}

Because Flirds and the in-run Shapley are values of the same game in the same
units, we can compare their \emph{magnitudes} directly. This comparison adds
information that Spearman's $\rho$ and Pearson's $r$ do not provide.
Table~\ref{tab:value-level} reports that comparison under full participation
($N{=}10$) using two scale-sensitive quantities. We evaluate only Flirds and
Flirds-1st because they alone produce values of the same game in the same
units as the reference. The relative $L_2$ error
$\|\hat\phi - \phi^{\mathrm{in}}\|_2 / \|\phi^{\mathrm{in}}\|_2$ is the
root-mean-square per-client error as a fraction of the target's own
root-mean-square size. Sign agreement is the fraction of clients on which the
two values agree in sign, which is the only property a sign gate reads.
Both are invariant to the global sign convention. Smaller is better for the
relative $L_2$ error and larger is better for sign agreement.

Three conditions behave as the correlations suggest and one does not. Under
the zero-update free-rider and label-flip, the relative $L_2$ error of Flirds
stays small in every partition. It grows in the clean condition, and under
gradient noise it grows further still, even though
Table~\ref{tab:inrun-fidelity} reports high $\rho$ for the two CIFAR-10
partitions in that condition. This is the regime
the bound of Appendix~\ref{app:remainder} anticipates, since
injected noise inflates $\|\Delta_S^r\|$ and the remainder grows with it.
Appendix~\ref{app:remainder-measured} measures the remainder in
this condition directly. We therefore read the high correlations in that
condition as evidence about ranking alone.

The two variants also separate on sign agreement. Flirds is perfect in 13 of
the 16 rows and remains close in the other three. Under gradient noise,
Flirds-1st drops on both CIFAR-10 partitions and assigns the wrong sign to four
of ten clients. These are precisely the errors on which a sign-based exclusion
rule acts.




\subsection{Agreement with the Retraining-Based Shapley}
\label{app:cnn-grid}

Table~\ref{tab:grid-retrain} scores six of the seven methods under full
participation ($N{=}10$) against the retraining-based Shapley, the common
external reference.
ComFedSV is omitted because these settings use full participation
(Appendix~\ref{app:train}), and both Shapley values are computed exhaustively:
the in-run Shapley over the $2^{10}$ coalitions of each round and the
retraining-based Shapley over the $2^{10}$ client subsets. The
CIFAR-10 Dirichlet rows are omitted because
Table~1 of the main paper already
reports them. The first column is not a method. It reports agreement between
the in-run and retraining-based Shapley values. That agreement varies across
settings. Its four-condition mean is highest for MNIST Dirichlet and lowest
for CIFAR-10 IID, whose clean condition reverses sign. For MNIST, the two games
align more closely under the non-IID partition than under its IID counterpart.
Flirds tracks this agreement closely in all three settings: its four-condition
mean stays within $0.035$ of the reference. Its discrepancy from the
retraining-based Shapley is therefore comparable to the discrepancy between
the two games. The bold entry marks the largest method mean, but, as in
Table~1 of the main paper, does not establish which method is best. The
retraining game renormalizes weights by subset size, and no evaluated method
targets it directly.

\begin{table}[!t]
\centering
\small
\setlength{\tabcolsep}{4pt}
% 추정치 두 열은 다른 표와 같은 순서(ablation -> full method)로 두고, 음영은
% 부록 C 앞머리의 Table conventions 대로 Flirds 열에 준다.
% 첫 열 헤더는 "Scale, partition, or condition" 이었을 때 그 열만 102.2pt 라 표
% 폭이 261.2pt 로 \columnwidth(237.6pt) 를 23.6pt 넘겨 두 줄로 나눠 두었다.
% partition 행이 없어(행은 model scale 아니면 condition 이고 partition 은 패널
% 헤더가 진다) "Scale or condition" 으로 줄이면 약 60pt 라 한 줄로도 표 폭이
% 224.8pt(블록 헤더 행이 결정) 안에 들어온다. 행 레이블 헤더를 둘째 줄에 두는
% 것은 tab:value-level / tab:taylor 와 같은 배치다.
\begin{tabular}{lcccB}
\hline
 & \multicolumn{2}{c}{in-run SV} & \multicolumn{2}{c}{estimator} \\
\cline{2-3}\cline{4-5}
Scale or condition & mean $\rho$ & min $\rho$ & Flirds-1st & Flirds \\
\hline
\noalign{\vskip\ourpanelsep}
\multicolumn{5}{l}{\emph{Alpaca, IID ($N{=}5$, full participation, $R{=}30$, clean)}} \\[\ourpanelsep]
1B & $-0.367$ & $-0.900$ & $-0.367$ & $-0.367$ \\
3B & $0.033$ & $-0.100$ & $0.033$ & $0.033$ \\
7B & $0.733$ & $0.500$ & $0.733$ & $0.733$ \\
\hline
\noalign{\vskip\ourpanelsep}
\multicolumn{5}{l}{\emph{Alpaca, IID ($N{=}20$, 2/20, $R{=}200$, clean)}} \\[\ourpanelsep]
1B & $-0.114$ & $-0.192$ & $-0.107$ & $-0.114$ \\
3B & $-0.243$ & $-0.447$ & $-0.227$ & $-0.243$ \\
7B & $0.164$ & $-0.036$ & $0.163$ & $0.172$ \\
\hline
\noalign{\vskip\ourpanelsep}
\multicolumn{5}{l}{\emph{GSM8K, IID ($N{=}50$, 5/50, $R{=}200$, 1B)}} \\[\ourpanelsep]
clean & $-0.026$ & $-0.177$ & $-0.025$ & $-0.025$ \\
answer-swap & $0.696$ & $0.658$ & $0.700$ & $0.699$ \\
zero-update & $0.748$ & $0.716$ & $0.748$ & $0.747$ \\
\hline
\noalign{\vskip\ourpanelsep}
\multicolumn{5}{l}{\emph{Five domains, non-IID ($N{=}5$, full participation, $R{=}10$, 1B)}} \\[\ourpanelsep]
clean & $0.867$ & $0.800$ & $0.867$ & $0.867$ \\
answer-swap & $0.933$ & $0.900$ & $0.933$ & $0.933$ \\
zero-update & $0.933$ & $0.900$ & $0.933$ & $0.933$ \\
\hline
\noalign{\vskip\ourpanelsep}
\multicolumn{5}{l}{\emph{Alpaca, IID ($N{=}5$, full participation, $R{=}10$, 1B)}} \\[\ourpanelsep]
clean & $0.133$ & $-0.300$ & $0.133$ & $0.133$ \\
\hline
\end{tabular}
\caption{Cross-seed Spearman's $\rho$ on the LLM track, over the three seed
pairs formed by $s \in \{0, 1, 2\}$ (settings as in
Appendices~\ref{app:train} and~\ref{app:data}). The \emph{in-run SV} columns give
the mean and the minimum over the three pairs for the reference value itself.
The last two columns give the mean for the two estimators of it.}
\label{tab:target-stability}
\end{table}

\begin{table*}[!t]
\centering
\small
\setlength{\tabcolsep}{4pt}
% 방법을 행, 조건을 열로 둔다(다른 표와 반대 방향). 원래 배치(조건이 행, 방법이 열)는
% ± 를 단 칸이 9열이라 내용만 485.7pt 여서 \textwidth(504pt) 안에 열 간격으로 쓸 수
% 있는 여유가 18pt 뿐이었고, 실제 표 폭이 517.7pt 로 13.7pt 넘쳤다(음영 열은 15.7pt).
% 전치하면 데이터 열이 4개로 줄어 폭이 약 314pt 가 되고 다른 표와 같은 4pt 간격을
% 쓸 수 있다. 값과 굵게 표기는 그대로 옮겼다. Flirds 음영은 B 열이 사라졌으므로
% \ourrow 행으로 주며, 이는 부록 C 앞머리 Table conventions 의 "shaded row or
% column" 규약 안에 있다(tab:gate-pr, tab:removal-cnn 과 같은 방식).
% 두 무대의 R 은 열 헤더에서 빠졌으므로 캡션이 어느 쪽이 R=10/30 인지 밝힌다.
% [S-LLM-A] LLM 트랙 retraining-based Shapley (부록 C.4). 출처 rundir:
%  - 5-도메인: runs/phase2_matrix/rundirs/1B_silo5_{clean,noisy,frzero}_aonly_s{0,1,2}
%    → runs/phase2_matrix/silo5_a_fidelity_1B.csv (merge_silo5_a.py 재생성),
%    열 spearman_a / pearson_a. in-run SV 행 = 그 파일의 (b)oracle 행.
%  - 앵커: runs/track_d/rundirs/1B_anchor5_seed{0,1,2}/phi.parquet
%    → runs/track_d/fidelity.csv (make_fidelity.py 재생성), 5-도메인과 같은 규약으로
%    열 spearman_a / pearson_a. in-run SV 행 = 그 파일의 (a)oracle 행(= (a) vs (b)).
% 3-seed mean ± 표본표준편차(ddof=1) — 본문 서술의 ±0.058/±0.115/±0.066 과 같은 규약
% (research-wiki 표는 ddof=0 이라 ±0.047 등으로 다르게 적힌다. 논문은 ddof=1 로 통일).
% ComFedSV 제외 사유는 tab:grid-retrain 과 동일(두 무대 모두 full participation).
\begin{tabular}{lcccc}
\hline
 & \multicolumn{3}{c}{Five-domain non-IID} & Alpaca, IID \\
\cline{2-4}\cline{5-5}
 & clean & answer-swap & zero-update & clean \\
\hline
\noalign{\vskip\ourpanelsep}
\multicolumn{5}{l}{\emph{in-run SV}} \\[\ourpanelsep]
rank $\rho$ & \ms{1.000}{0.000} & \ms{1.000}{0.000} & \ms{1.000}{0.000} & \ms{0.933}{0.058} \\
value $r$ & \ms{1.000}{0.000} & \ms{1.000}{0.000} & \ms{1.000}{0.000} & \ms{0.933}{0.066} \\
\hline
\noalign{\vskip\ourpanelsep}
\multicolumn{5}{l}{\emph{Spearman's $\rho$ against the retraining-based Shapley}} \\[\ourpanelsep]
GTG-Shapley & \msb{1.000}{0.000} & \msb{1.000}{0.000} & \msb{1.000}{0.000} & \msb{0.933}{0.058} \\
FedSV & \ms{0.933}{0.058} & \msb{1.000}{0.000} & \ms{0.933}{0.058} & \ms{0.733}{0.208} \\
ShapleyFL & \msb{1.000}{0.000} & \msb{1.000}{0.000} & \msb{1.000}{0.000} & \ms{0.767}{0.404} \\
FedIF & \ms{0.867}{0.115} & \ms{0.933}{0.058} & \ms{0.900}{0.100} & \ms{0.167}{0.751} \\
Flirds-1st & \msb{1.000}{0.000} & \msb{1.000}{0.000} & \msb{1.000}{0.000} & \msb{0.933}{0.058} \\
\ourrow Flirds & \msb{1.000}{0.000} & \msb{1.000}{0.000} & \msb{1.000}{0.000} & \msb{0.933}{0.058} \\
\hline
\end{tabular}
\caption{Spearman's $\rho$ against the retraining-based Shapley on the LLM
track (five domains and Alpaca, Llama-3.2-1B-Instruct; $N{=}5$, full
participation, $R{=}10$ for the five-domain corpus and $R{=}30$ for Alpaca).
The two \emph{in-run SV} rows give the agreement between the retraining-based
and the in-run Shapley, in rank ($\rho$) and in value ($r$).}
\label{tab:llm-retrain}
\end{table*}

\begin{table*}[t]
\centering
\small
\setlength{\tabcolsep}{4pt}
% [T1] Taylor remainder 직접 측정. 출처 rundir: runs/taylor_remainder/rundirs/
% {1B_silo5_<threat>_taylor_seed{0,1,2}, cifar10_dir1_<threat>_taylor_seed{0,1,2}}
% /summary.json (pooled.resid{1,2}.{median,max} / mean_abs_u_grand /
% frac_t2_le_t1 / resid2_median_over_ulp / rounds[].norm_dW). 3-seed 평균.
% 표 전체를 python runs/taylor_remainder/make_analysis.py 로 재생성한다.
% order(loglog / scale-sweep 기울기)는 싣지 않는다 — 이 절의 의무는 magnitude 이고
% 차수는 Appendix~\ref{app:remainder} 의 정리가 담당한다.
% "first & second" 가 두 지표에 대해 반복되므로 \cline 으로 묶는다(tab:value-level
% 과 같은 규약).
\begin{tabular}{lcccccc}
\hline
 & \multicolumn{2}{c}{median remainder} & \multicolumn{2}{c}{median\,/\,$|u_r(P_r)|$} & & \\
\cline{2-3}\cline{4-5}
Condition & first & second & first & second & 2nd $\le$ 1st & 2nd\,/\,floor \\
\hline
\noalign{\vskip\ourpanelsep}
\multicolumn{7}{l}{\emph{Five-domain non-IID, Llama-3.2-1B-Instruct ($N{=}5$, full participation, $R{=}10$)}} \\[\ourpanelsep]
clean & $1.4\times10^{-6}$ & \vb{4.8\times10^{-7}} & $0.09\%$ & \vb{0.03\%} & $83.1\%$ & $2.0$ \\
answer-swap & $1.4\times10^{-6}$ & \vb{5.0\times10^{-7}} & $0.10\%$ & \vb{0.03\%} & $81.1\%$ & $2.1$ \\
zero-update & $1.2\times10^{-6}$ & \vb{3.6\times10^{-7}} & $0.10\%$ & \vb{0.03\%} & $82.4\%$ & $1.5$ \\
\hline
\noalign{\vskip\ourpanelsep}
\multicolumn{7}{l}{\emph{CIFAR-10 / Dir(1), FedSVCNN ($N{=}10$, full participation, $R{=}10$)}} \\[\ourpanelsep]
clean & $2.2\times10^{-2}$ & \vb{1.6\times10^{-3}} & $17.4\%$ & \vb{1.3\%} & $97.1\%$ & $1.2\times10^{4}$ \\
zero-update & $8.4\times10^{-3}$ & \vb{9.1\times10^{-4}} & $8.6\%$ & \vb{0.9\%} & $99.2\%$ & $6.4\times10^{3}$ \\
gradient noise & $3.8\times10^{-1}$ & \vb{9.9\times10^{-2}} & $271\%$ & \vb{70.4\%} & $94.8\%$ & $6.5\times10^{5}$ \\
label-flip & $1.2\times10^{-2}$ & \vb{2.5\times10^{-3}} & $12.8\%$ & \vb{2.6\%} & $96.8\%$ & $1.7\times10^{4}$ \\
\hline
\end{tabular}
\caption{Taylor remainder of the round utility, over all $2^{|P_r|}$ coalitions
of every round. $|u_r(P_r)|$ is the mean per-round realized validation-loss
decrease. ``2nd $\le$ 1st'' is the fraction of enumerated coalitions on which
the second-order surrogate is at least as accurate as the first. ``floor'' is
one unit in the last place of the round's base loss in fp32. All quantities
are averaged over three seeds.}
\label{tab:taylor}
\end{table*}


\subsection{Retraining-Based Shapley on the LLM Track}
\label{app:llm-fidelity}

We compute the retraining-based Shapley in two LLM settings
(Table~\ref{tab:llm-retrain}). In the five-domain non-IID setting
($N{=}5$, full participation, $R{=}10$), we exhaustively enumerate the
classical value through $2^5$ retraining runs from scratch and compare it with
in-run enumeration on the corresponding realized run. The two Shapley values
agree perfectly in rank and value under all three conditions: clean,
answer-swap, and zero-update. Consequently, each method has the same agreement
with both references. GTG-Shapley, ShapleyFL, and both Flirds variants
reproduce the ranking exactly in all three conditions. FedSV does not under
clean and zero-update, and FedIF does not in any condition. This setting is a
measured case in which the classical, method-neutral value directly supports
the in-run reference used throughout the paper.

In the Alpaca setting at $N{=}5$ ($R{=}30$), the two Shapley values agree
closely but not exactly. We therefore compare every method directly with the
retraining-based Shapley. GTG-Shapley and both Flirds variants reproduce the
in-run ranking exactly in every seed, so their direct agreement equals the
agreement between the two reference values. FedSV, ShapleyFL, and FedIF do not
reproduce that ranking and fall below this level, with FedIF farthest below.
Because both settings contain only five clients, their rankings have limited
room to differ and their seed spreads are wide. We therefore treat the
ordering as indicative rather than as a separation among methods.
Appendix~\ref{app:signal} reports the cross-seed reproducibility of both
settings. The Alpaca setting is IID and clean, and the cross-seed
reproducibility of its in-run Shapley is negative
(Table~\ref{tab:target-stability}), whereas the five-domain non-IID setting is
high.

\subsection{Direct Measurement of the Taylor Remainder}
\label{app:remainder-measured}

The preceding subsections compare contribution values, either with their
direct target or with the common external reference.
This subsection instead measures the per-round Taylor
remainder underlying the difference between
Flirds and its target,
$|u_r(S) - \hat u_r(S)|$ from Appendix~\ref{app:remainder}.

\noindent\textbf{Protocol.} We measure the remainder in two settings, one on each track,
chosen so that both admit exhaustive enumeration of every coalition: the five-domain
non-IID LLM setting (Llama-3.2-1B-Instruct, $N{=}5$, full participation,
$R{=}10$; Appendices~\ref{app:train} and~\ref{app:data}), and the
CNN setting of Table~1 of the main paper (FedSVCNN on CIFAR-10 /
Dirichlet($\alpha{=}1$), $N{=}10$, full participation, $R{=}10$). The CNN
setting is included because the bound of
Appendix~\ref{app:remainder} is conditional on a $C^3$ assumption that a ReLU
network with max-pooling does not satisfy, so on that track a direct
measurement is the only evidence available. Each setting is measured under
every condition
reported for its track in Section~5.1 of the main paper. For every one of the
$2^{|P_r|}$ coalitions in each round we evaluate the round utility $u_r(S)$ of
Eq.~(5) of the main paper together with its first- and second-order surrogates
and record the absolute differences $|u_r(S) - \hat u_r^{(1)}(S)|$ and
$|u_r(S) - \hat u_r(S)|$. Remainders are pooled within a seed and averaged over
three seeds. We also report the fp32 evaluation floor, one unit in the last
place of the round's base loss, so that the resolution of each measurement can
be read directly.

\noindent\textbf{Result.} Table~\ref{tab:taylor} reports both remainders for all seven
setting--condition pairs. The second-order median remainder is smaller than the
first-order median in every pair. Outside gradient noise, its median is $0.03\%$ of
the mean per-round realized validation-loss decrease $|u_r(P_r)|$ on the LLM
track and $0.9\%$--$2.6\%$ on the CNN track. Retaining the curvature term
reduces the remainder by $2.8$--$3.3\times$ on the LLM track and
$4.8$--$13.8\times$ on the CNN track. The second-order surrogate is at least
as accurate as the first outside gradient noise on $81$--$83\%$ and
$97$--$99\%$ of the enumerated coalitions, respectively. Across the four CNN
conditions, the first-order
ratio also tracks the player-level relative $L_2$ error reported for
Flirds-1st in Table~\ref{tab:value-level}, the quantity into which the
remainder propagates.

The two tracks differ in measurement resolution. On the LLM track, the
second-order remainder is only $1.5$--$2.1$ times one unit in the last place of
the base loss. The reported value is therefore an upper bound. The
remainder may be smaller. On the CNN track, it exceeds the evaluation floor by
more than three orders of magnitude in every condition and is fully resolved.
Thus, both the floor-limited LLM measurement and the resolved CNN measurement
place the median remainder at a small fraction of the quantity being attributed
outside gradient noise.

\noindent\textbf{The gradient-noise condition.} Additive gradient noise departs from
this pattern. In Table~\ref{tab:taylor}, the first-order median remainder
exceeds the quantity being attributed, while the smaller second-order remainder
reaches a fraction of it.
This result is consistent with the observation in Section~4.3 of the main paper
that a round displacement need not be small and the remainder is not
automatically negligible. Appendix~\ref{app:value-level} accordingly
interprets the correlations in this condition as evidence about ranking
alone. The measurement here verifies the remainder underlying that
interpretation. The round-level ratios and the player-level relative $L_2$
errors in Table~\ref{tab:value-level} use different normalizers,
$|u_r(P_r)|$ and $\|\phi^{\mathrm{in}}\|_2$, but capture the same effect at two
levels and agree in size at both orders.

The enlarged remainder does not disturb the decision-relevant ranking and
signs. Flirds keeps perfect sign agreement in this condition, whereas
Flirds-1st assigns the wrong sign to four of ten clients
(Table~\ref{tab:value-level}). The two variants separate in the same way
downstream. Flirds-1st alone reverses rank agreement in
Table~1 of the main paper, after maintaining it in every other
condition, and recovers none of the
accuracy lost to corruption in Table~\ref{tab:online}. Flirds maintains its
rank agreement and recovers most of that accuracy. The curvature term
separates the variants in both results.



\begin{table*}[!t]
\centering
\small
% mean 열은 조건 열들의 집계이므로 세로선으로 분리하고, 굵게는 mean 열에만 쓴다.
\begin{tabular}{lcccc|c}
\hline
 & clean & zero-update & gradient noise & label-flip & mean \\
\hline
\noalign{\vskip\ourpanelsep}
\multicolumn{6}{l}{\emph{CIFAR-10 / Dir(1), online}} \\[\ourpanelsep]
vanilla & \ms{63.89}{0.52} & \ms{58.79}{0.29} & \ms{24.36}{2.22} & \ms{52.47}{2.89} & $49.88$ \\
oracle-exclusion & --- & \ms{62.03}{0.28} & \ms{62.03}{0.28} & \ms{62.36}{0.31} & --- \\
selection-random & --- & \ms{58.38}{2.02} & \ms{25.90}{2.08} & \ms{50.18}{6.09} & --- \\
\hline
GTG-Shapley & \ms{60.51}{0.74} & \ms{39.15}{0.81} & \ms{59.72}{0.52} & \ms{54.79}{1.15} & $53.54$ \\
FedSV & \ms{59.82}{0.77} & \ms{39.66}{1.13} & \ms{59.72}{0.61} & \ms{51.64}{0.83} & $52.71$ \\
ComFedSV & \ms{59.63}{0.59} & \ms{39.18}{1.04} & \ms{58.71}{0.54} & \ms{51.52}{2.78} & $52.26$ \\
ShapleyFL & \ms{60.45}{0.52} & \ms{40.20}{2.14} & \ms{61.15}{1.04} & \ms{52.78}{1.65} & $53.65$ \\
FedIF & \ms{63.86}{0.10} & \ms{61.43}{1.17} & \ms{24.79}{1.32} & \ms{57.28}{0.87} & $51.84$ \\
Flirds-1st & \ms{63.84}{0.17} & \ms{62.16}{0.38} & \ms{24.79}{1.32} & \ms{57.17}{0.81} & $51.99$ \\
\ourrow Flirds & \ms{63.15}{0.74} & \ms{61.48}{0.02} & \ms{56.68}{2.60} & \ms{57.12}{0.84} & \vb{59.61} \\
\hline
\noalign{\vskip\ourpanelsep}
\multicolumn{6}{l}{\emph{CIFAR-10 / IID, online}} \\[\ourpanelsep]
vanilla & \ms{64.88}{0.05} & \ms{60.83}{0.46} & \ms{25.64}{0.40} & \ms{51.71}{1.84} & $50.77$ \\
oracle-exclusion & --- & \ms{63.56}{0.11} & \ms{63.56}{0.11} & \ms{63.10}{0.39} & --- \\
selection-random & --- & \ms{59.86}{0.86} & \ms{26.45}{1.05} & \ms{50.23}{5.36} & --- \\
\hline
GTG-Shapley & \ms{63.86}{0.32} & \ms{55.53}{1.39} & \ms{61.56}{0.25} & \ms{61.72}{0.59} & $60.67$ \\
FedSV & \ms{61.98}{0.29} & \ms{51.25}{0.47} & \ms{62.07}{0.89} & \ms{62.24}{0.18} & $59.39$ \\
ComFedSV & \ms{61.98}{0.29} & \ms{51.25}{0.47} & \ms{62.07}{0.89} & \ms{62.24}{0.18} & $59.39$ \\
ShapleyFL & \ms{62.39}{0.38} & \ms{51.03}{0.88} & \ms{61.85}{0.78} & \ms{62.64}{0.33} & $59.48$ \\
FedIF & \ms{64.84}{0.21} & \ms{63.42}{0.08} & \ms{26.19}{0.73} & \ms{59.76}{0.19} & $53.55$ \\
Flirds-1st & \ms{64.84}{0.21} & \ms{63.42}{0.08} & \ms{26.19}{0.73} & \ms{59.70}{2.66} & $53.54$ \\
\ourrow Flirds & \ms{64.28}{0.76} & \ms{63.08}{0.50} & \ms{61.43}{2.11} & \ms{59.67}{0.77} & \vb{62.12} \\
\hline
\noalign{\vskip\ourpanelsep}
\multicolumn{6}{l}{\emph{CIFAR-10 / IID, selection-retrain}} \\[\ourpanelsep]
vanilla & \ms{64.88}{0.05} & \ms{60.83}{0.46} & \ms{25.64}{0.40} & \ms{51.71}{1.84} & $50.77$ \\
oracle-exclusion & --- & \ms{63.56}{0.11} & \ms{63.56}{0.11} & \ms{63.10}{0.39} & --- \\
selection-random & --- & \ms{59.86}{0.86} & \ms{26.45}{1.05} & \ms{50.23}{5.36} & --- \\
\hline
GTG-Shapley & \ms{63.70}{0.27} & \ms{57.83}{1.06} & \ms{63.52}{0.11} & \ms{63.05}{0.28} & $62.03$ \\
FedSV & \ms{63.75}{0.14} & \ms{49.27}{1.41} & \ms{63.52}{0.11} & \ms{63.05}{0.28} & $59.90$ \\
ComFedSV & \ms{63.75}{0.14} & \ms{49.27}{1.41} & \ms{63.52}{0.11} & \ms{63.05}{0.28} & $59.90$ \\
ShapleyFL & \ms{63.49}{0.08} & \ms{47.63}{2.60} & \ms{63.52}{0.11} & \ms{63.05}{0.28} & $59.42$ \\
FedIF & \ms{64.56}{0.43} & \ms{63.52}{0.11} & \ms{25.64}{0.40} & \ms{63.05}{0.28} & $54.19$ \\
Flirds-1st & \ms{64.56}{0.43} & \ms{63.52}{0.11} & \ms{25.64}{0.40} & \ms{63.05}{0.28} & $54.19$ \\
\ourrow Flirds & \ms{64.19}{0.46} & \ms{63.26}{0.16} & \ms{62.57}{0.23} & \ms{63.05}{0.28} & \vb{63.27} \\
\hline
\end{tabular}
\caption{Test accuracy (\%) on the CNN track under online sign-gating and
under sign-gating with selection-retrain (CIFAR-10; $N{=}100$, 10/100
participation, $R{=}120$).}
\label{tab:online}
\end{table*}

% ===================================================================== D
\section{Reproducibility of the Target Signal}
\label{app:signal}

\begin{table}[!t]
\centering
\small
\setlength{\tabcolsep}{6pt}
\begin{tabular}{lcc}
\hline
 & LLM & CNN \\
 & (13 / 10 entries) & (16 entries) \\
\hline
GTG-Shapley & $0.038$ & $0.171$ \\
FedSV & $0.153$ & $0.145$ \\
ShapleyFL & $0.121$ & $0.133$ \\
FedIF & $0.195$ & $0.141$ \\
Flirds-1st & $0.002$ & $0.144$ \\
\ourrow Flirds & \vb{0.001} & \vb{0.050} \\
\hline
\end{tabular}
\caption{Mean absolute difference between a method's cross-seed agreement and
that of the in-run Shapley it estimates. Smaller is closer to the target's own
reproducibility. The LLM column covers the rows of
Table~\ref{tab:target-stability}: thirteen for the two Flirds variants and the
ten outside the GSM8K setting for the other methods, which that setting does
not evaluate. The CNN column covers the sixteen setting--condition
combinations of CIFAR-10 and MNIST
under full participation ($N{=}10$, $R{=}10$, both partitions, four
conditions).}
\label{tab:target-inherit}
\end{table}

The clean condition behaves differently from the corruption conditions in
Section~5.2 of the main paper and Appendix~\ref{app:fidelity}: its correlations
are lower and vary more widely across seeds. The results below show that this
difference lies in the reference value rather than in the estimators evaluated
against it. A correlation with a reference ranking transfers beyond its
realized run only when that ranking is itself a stable property of the clients
\citep{geimer2025volatility}. When clients do not differ systematically, their
ranking changes across seeds.

Table~\ref{tab:target-stability} measures this stability by comparing the
in-run Shapley between seed pairs. Its first two columns report the mean and
minimum Spearman's $\rho$ over the three pairs. Larger values indicate a more
reproducible ranking. When cross-seed agreement is low, a high within-run
correlation can still accurately reproduce a ranking that is itself unstable.

The measurement spans a wide range. Agreement is highest in the five-domain
setting with one domain per client and remains high under every condition. It
is lowest in the 1B and 3B Alpaca IID settings, where three of the four values
are negative. It is also near zero in the clean condition of the main LLM
setting, which compares 50 clients per seed pair rather than five. In both
settings that vary model scale, agreement is higher at 7B than at 1B. Under
full participation, it is high at 7B. The settings also differ in corpus,
partition, and the configuration associated with model scale
(Appendices~\ref{app:train} and~\ref{app:data}). The table therefore identifies
where the reference ranking is reproducible, not what causes its
reproducibility.

The last two columns of Table~\ref{tab:target-stability} apply the same
measurement to the two estimators. Both reproduce the cross-seed profile of
the in-run Shapley as closely as they reproduce its within-run values, staying
level with the reference across all thirteen rows.
Table~\ref{tab:target-inherit} extends this comparison
to the other methods and to the CNN track under full participation. ComFedSV
is omitted because only three of the twenty-six combinations it would be
scored on use partial participation (Appendix~\ref{app:cnn-grid}). A method
that reproduces its target in every seed must also reproduce the target's
agreement across seeds, as expected from the fidelity in
Section~5.2 of the main paper. Methods that do not reproduce the target also
depart more from its
cross-seed profile. On the LLM track, the two Flirds variants stay an order of
magnitude closer to the reference than every other method. On the CNN track,
Flirds again stays closest, and every other method, including Flirds-1st,
departs more than twice as far.

Table~\ref{tab:target-inherit} reports the absolute difference from the target's
cross-seed agreement, not the direction of each departure. It therefore
measures whether a method inherits the target's reproducibility profile rather
than whether the method's own cross-seed agreement is high. Cross-seed
agreement and within-run fidelity must be read together: a departure from the
target's profile does not by itself identify a better estimate of the round
game.

The CIFAR-10 IID clean condition provides two displayed signs of a weak
ordering signal. The
agreement between the two Shapley values is negative there
(Table~\ref{tab:grid-retrain}), and the clean row of the removal curve puts
every method near zero, so the removal order does not change performance
(Table~\ref{tab:removal-cnn}). Both describe a setting in which client
order carries little information, rather than one in which the methods fail.

Cross-seed agreement and within-run fidelity support two readings of a
contribution value. Within-run fidelity asks how closely an estimate
reproduces the in-run Shapley of one realized run. Section~5.2 of the main paper
and Appendix~\ref{app:fidelity} measure this quantity, and
Section~5.4 of the main paper
compares its computational cost. This reading does not require stability
across seeds because it concerns one trajectory. Cross-seed agreement instead
asks whether the induced ordering is a durable client property
(Table~\ref{tab:target-stability}). When that agreement is low, the second
reading is unavailable to every method in that setting.
Table~\ref{tab:target-inherit} places this attenuation in the target, which
Flirds inherits, rather than in the estimator.

The attenuation is therefore a property of the setting, not of one estimator.
Within a given setting, condition, and seed, every method is evaluated against
the same reference on the same realized trajectory, so the effect is shared. We therefore
report and aggregate the clean condition on the same footing as the others and
interpret it as the control: it records what the methods do when little
separates the clients. This also motivates reporting cross-seed reproducibility
alongside within-run fidelity.


\section{Full Intervention Results}
\label{app:intervention}

The sign-gating rule of Section~5.3 of the main paper is applied in two ways.
Under \emph{selection-retrain}, contributions accumulate over a complete
initial run. Clients with accumulated contribution at most zero are then
dropped, and training restarts from the same initial weights using only the
retained clients. The gate therefore fires once using evidence from the whole
run. Under \emph{online} gating, the same rule operates during a single
training run. Appendix~\ref{app:online} gives its three additional parameters.
Section~5.3 of the main paper reports selection-retrain for the CIFAR-10
Dirichlet partition and the main LLM setting.
This appendix reports online gating, the MNIST and IID
repetitions, the clients excluded by the LLM gate, and a removal curve that
evaluates the contribution ranking through retraining performance.

The vanilla row in the tables is the initial run that provides the
contribution estimates. When the gate excludes no clients, retraining is
identical to training without intervention, so we report the corresponding
result from that run. More generally, a retrained run is determined by the
exclusions and the seed alone, so any two rows whose exclusions agree in
every seed report the same mean and the same spread.

\begin{table*}[!t]
\centering
\small
% mean 열은 조건 열들의 집계이므로 세로선으로 분리하고, 굵게는 mean 열에만 쓴다.
\begin{tabular}{lcccc|c}
\hline
 & clean & zero-update & gradient noise & label-flip & mean \\
\hline
\noalign{\vskip\ourpanelsep}
\multicolumn{6}{l}{\emph{MNIST / Dir(1), online}} \\[\ourpanelsep]
vanilla & \ms{97.92}{0.21} & \ms{97.13}{0.21} & \ms{89.64}{1.37} & \ms{96.25}{0.65} & $95.24$ \\
oracle-exclusion & --- & \ms{97.76}{0.21} & \ms{97.76}{0.21} & \ms{97.87}{0.30} & --- \\
selection-random & --- & \ms{97.11}{0.21} & \ms{90.42}{1.49} & \ms{96.54}{1.00} & --- \\
\hline
GTG-Shapley & \ms{97.85}{0.17} & \ms{95.16}{0.47} & \ms{97.82}{0.08} & \ms{97.81}{0.30} & $97.16$ \\
FedSV & \ms{97.84}{0.26} & \ms{93.44}{0.40} & \ms{97.75}{0.13} & \ms{97.73}{0.26} & $96.69$ \\
ComFedSV & \ms{97.87}{0.18} & \ms{93.19}{0.45} & \ms{97.80}{0.10} & \ms{97.70}{0.30} & $96.64$ \\
ShapleyFL & \ms{97.78}{0.20} & \ms{92.75}{0.29} & \ms{97.96}{0.25} & \ms{97.80}{0.32} & $96.57$ \\
FedIF & \ms{97.92}{0.31} & \ms{97.84}{0.14} & \ms{90.95}{0.21} & \ms{97.61}{0.19} & $96.08$ \\
Flirds-1st & \ms{97.89}{0.28} & \ms{97.89}{0.21} & \ms{90.91}{0.72} & \ms{97.74}{0.33} & $96.11$ \\
\ourrow Flirds & \ms{97.92}{0.26} & \ms{97.27}{0.85} & \ms{97.72}{0.22} & \ms{97.74}{0.54} & \vb{97.66} \\
\hline
\noalign{\vskip\ourpanelsep}
\multicolumn{6}{l}{\emph{MNIST / Dir(1), selection-retrain}} \\[\ourpanelsep]
vanilla & \ms{97.92}{0.21} & \ms{97.13}{0.21} & \ms{89.64}{1.37} & \ms{96.25}{0.65} & $95.24$ \\
oracle-exclusion & --- & \ms{97.76}{0.21} & \ms{97.76}{0.21} & \ms{97.87}{0.30} & --- \\
selection-random & --- & \ms{97.11}{0.21} & \ms{90.42}{1.49} & \ms{96.54}{1.00} & --- \\
\hline
GTG-Shapley & \ms{97.99}{0.07} & \ms{97.14}{0.25} & \ms{97.78}{0.21} & \ms{97.89}{0.34} & $97.70$ \\
FedSV & \ms{97.98}{0.14} & \ms{96.78}{0.41} & \ms{97.79}{0.17} & \ms{97.89}{0.34} & $97.61$ \\
ComFedSV & \ms{98.01}{0.13} & \ms{96.79}{0.42} & \ms{97.79}{0.17} & \ms{97.89}{0.34} & $97.62$ \\
ShapleyFL & \ms{97.97}{0.12} & \ms{96.77}{0.43} & \ms{97.79}{0.17} & \ms{97.89}{0.34} & $97.61$ \\
FedIF & \ms{97.93}{0.15} & \ms{97.83}{0.18} & \ms{89.64}{1.37} & \ms{97.87}{0.30} & $95.82$ \\
Flirds-1st & \ms{97.95}{0.19} & \ms{97.82}{0.15} & \ms{89.64}{1.37} & \ms{97.87}{0.30} & $95.82$ \\
\ourrow Flirds & \ms{98.10}{0.23} & \ms{97.84}{0.21} & \ms{97.55}{0.36} & \ms{97.90}{0.32} & \vb{97.85} \\
\hline
\noalign{\vskip\ourpanelsep}
\multicolumn{6}{l}{\emph{MNIST / IID, online}} \\[\ourpanelsep]
vanilla & \ms{98.16}{0.22} & \ms{97.41}{0.15} & \ms{92.21}{0.81} & \ms{96.31}{0.16} & $96.02$ \\
oracle-exclusion & --- & \ms{98.07}{0.09} & \ms{98.07}{0.09} & \ms{98.09}{0.17} & --- \\
selection-random & --- & \ms{97.45}{0.45} & \ms{91.57}{1.55} & \ms{96.79}{0.83} & --- \\
\hline
GTG-Shapley & \ms{98.13}{0.17} & \ms{97.11}{0.47} & \ms{98.00}{0.28} & \ms{97.95}{0.17} & $97.80$ \\
FedSV & \ms{98.09}{0.07} & \ms{96.42}{0.81} & \ms{98.03}{0.19} & \ms{98.07}{0.18} & $97.65$ \\
ComFedSV & \ms{98.09}{0.07} & \ms{96.42}{0.81} & \ms{98.03}{0.19} & \ms{98.07}{0.18} & $97.65$ \\
ShapleyFL & \ms{98.05}{0.19} & \ms{96.48}{0.52} & \ms{98.02}{0.20} & \ms{98.10}{0.16} & $97.66$ \\
FedIF & \ms{98.12}{0.20} & \ms{98.02}{0.09} & \ms{91.44}{1.22} & \ms{98.06}{0.18} & $96.41$ \\
Flirds-1st & \ms{98.16}{0.17} & \ms{97.77}{0.52} & \ms{91.44}{1.22} & \ms{98.09}{0.22} & $96.37$ \\
\ourrow Flirds & \ms{97.90}{0.13} & \ms{97.97}{0.22} & \ms{97.81}{0.25} & \ms{97.95}{0.31} & \vb{97.91} \\
\hline
\noalign{\vskip\ourpanelsep}
\multicolumn{6}{l}{\emph{MNIST / IID, selection-retrain}} \\[\ourpanelsep]
vanilla & \ms{98.16}{0.22} & \ms{97.41}{0.15} & \ms{92.21}{0.81} & \ms{96.31}{0.16} & $96.02$ \\
oracle-exclusion & --- & \ms{98.07}{0.09} & \ms{98.07}{0.09} & \ms{98.09}{0.17} & --- \\
selection-random & --- & \ms{97.45}{0.45} & \ms{91.57}{1.55} & \ms{96.79}{0.83} & --- \\
\hline
GTG-Shapley & \ms{98.14}{0.08} & \ms{97.68}{0.20} & \ms{98.07}{0.09} & \ms{98.09}{0.17} & $98.00$ \\
FedSV & \ms{98.09}{0.17} & \ms{97.42}{0.29} & \ms{98.07}{0.09} & \ms{98.09}{0.17} & $97.92$ \\
ComFedSV & \ms{98.09}{0.17} & \ms{97.42}{0.29} & \ms{98.07}{0.09} & \ms{98.09}{0.17} & $97.92$ \\
ShapleyFL & \ms{98.04}{0.15} & \ms{97.41}{0.15} & \ms{98.07}{0.09} & \ms{98.09}{0.17} & $97.90$ \\
FedIF & \ms{98.12}{0.13} & \ms{98.07}{0.09} & \ms{91.83}{1.46} & \ms{98.09}{0.17} & $96.53$ \\
Flirds-1st & \ms{98.12}{0.13} & \ms{98.07}{0.09} & \ms{91.83}{1.46} & \ms{98.09}{0.17} & $96.53$ \\
\ourrow Flirds & \ms{98.15}{0.13} & \ms{98.04}{0.12} & \ms{98.09}{0.16} & \ms{98.09}{0.17} & \vb{98.09} \\
\hline
\end{tabular}
\caption{Test accuracy (\%) on the CNN track under online sign-gating and
under sign-gating with selection-retrain (MNIST, both partitions; $N{=}100$,
10/100 participation, $R{=}120$).}
\label{tab:mnist-int}
\end{table*}

\begin{table*}[!t]
\centering
\small
\setlength{\tabcolsep}{4pt}
% 두 프로토콜은 행 레이블이 같으므로 스택 대신 열 그룹으로 나란히 둔다(패널 사이는
% tab:inrun-fidelity 와 같은 \qquad). 굵게는 두 오염 조건의 precision/recall 에만
% 쓴다: clean 에는 오염 클라이언트가 없어 두 지표가 비교 축이 아니고, false excl.
% 은 게이트가 덜 발화할수록 낮아지므로 낮은 값이 곧 우월이 아니다(부록 C 앞머리의
% Table conventions 중 "entries that are not part of a comparison carry no bold").
\begin{tabular}{llccc@{\qquad}ccc}
\hline
 & & \multicolumn{3}{c}{\emph{selection-retrain (per client)}} & \multicolumn{3}{c}{\emph{online (per client-round)}} \\
\cline{3-5}\cline{6-8}
 & Condition & precision & recall & false excl. & precision & recall & false excl. \\
\hline
Flirds-1st & clean & --- & --- & \ms{0.00}{0.00} & \ms{0.00}{0.00} & --- & \ms{1.60}{0.46} \\
Flirds-1st & answer-swap & \msb{100.00}{0.00} & \ms{66.67}{5.77} & \ms{0.00}{0.00} & \msb{85.69}{0.37} & \ms{4.56}{0.43} & \ms{0.51}{0.06} \\
Flirds-1st & zero-update & \msb{100.00}{0.00} & \msb{100.00}{0.00} & \ms{0.00}{0.00} & \msb{95.80}{0.51} & \msb{94.54}{0.68} & \ms{2.77}{0.34} \\
\hline
\ourrow Flirds & clean & --- & --- & \ms{0.00}{0.00} & \ms{0.00}{0.00} & --- & \ms{1.89}{0.43} \\
\ourrow Flirds & answer-swap & \msb{100.00}{0.00} & \msb{75.00}{10.00} & \ms{0.00}{0.00} & \ms{79.28}{2.82} & \msb{5.54}{0.28} & \ms{0.96}{0.12} \\
\ourrow Flirds & zero-update & \msb{100.00}{0.00} & \msb{100.00}{0.00} & \ms{0.00}{0.00} & \ms{95.16}{0.74} & \msb{94.54}{0.68} & \ms{3.21}{0.51} \\
\hline
\end{tabular}
\caption{Precision, recall, and false-exclusion rate (\%) of the gate under
sign-gating with selection-retrain and under online sign-gating (GSM8K,
Llama-3.2-1B-Instruct; $N{=}50$, 5/50 participation, $R{=}200$).}
\label{tab:gate-pr}
\end{table*}


\begin{table*}[!t]
\centering
\small
% mean 열은 조건 열들의 집계이므로 세로선으로 분리하고, 굵게는 mean 열에만 쓴다.
% 이 표의 mean 최고값은 우리 방법이 아니라 FedIF($+15.09$)다. in-run SV 행은
% reference 이므로 비교 대상에서 빠진다(부록 C 앞머리의 Table conventions).
\begin{tabular}{lcccc|c}
\hline
 & clean & zero-update & gradient noise & label-flip & mean \\
\hline
in-run SV & \ms{-0.27}{0.97} & \ms{+21.72}{0.34} & \ms{+24.57}{2.56} & \ms{+13.98}{0.40} & $+15.00$ \\
\hline
GTG-Shapley & \ms{+0.10}{0.82} & \ms{+21.68}{0.68} & \ms{+23.60}{2.34} & \ms{+14.09}{0.68} & $+14.87$ \\
FedSV & \ms{+0.14}{0.24} & \ms{+10.50}{10.06} & \ms{+23.78}{1.87} & \ms{+14.80}{1.21} & $+12.31$ \\
ShapleyFL & \ms{-0.20}{0.93} & \ms{-16.26}{9.45} & \ms{+24.62}{2.47} & \ms{+14.51}{0.64} & $+5.67$ \\
FedIF & \ms{-0.41}{0.30} & \ms{+21.95}{0.32} & \ms{+24.22}{1.99} & \ms{+14.61}{0.37} & \vb{+15.09} \\
Flirds-1st & \ms{+0.03}{0.35} & \ms{+21.89}{0.26} & \ms{-1.89}{4.76} & \ms{+14.31}{0.66} & $+8.59$ \\
\ourrow Flirds & \ms{-0.33}{0.69} & \ms{+21.76}{0.30} & \ms{+24.45}{2.31} & \ms{+13.84}{0.37} & $+14.93$ \\
\hline
\end{tabular}
\caption{Mean gap between the two removal curves, in percentage points
(CIFAR-10, IID; $N{=}10$, full participation, $R{=}10$, 4/10 corrupted).}
\label{tab:removal-cnn}
\end{table*}


\subsection{Online Gating}
\label{app:online}

The first two panels of Table~\ref{tab:online} apply the gate online in the
same setting and with the same contribution estimates as
Section~5.3 of the main paper. The rule that
decides exclusion is the one used there, dropping a client whose accumulated
contribution is at most zero. Because the rule now runs during training, it
requires three parameters shared by every method and never tuned per
setting or condition. Gating is disabled for the first 10 rounds, and a client is judged only after
participating in at least 2 rounds. Afterward, the excluded clients rotate back
into the candidate pool one at a time every fifth round, so exclusion is not absorbing
and a client's accumulated contribution can recover. Contributions accumulate over rounds without
decay. The ordering from selection-retrain is
broadly preserved, but differences narrow because gating is disabled for the
first 10 rounds and the resulting model already reflects updates from corrupted
clients during that interval. The first-order methods stay at vanilla's level
under gradient noise, whereas Flirds recovers, gaining more than $30$ points
over them. The renormalizing methods again fall below vanilla under the
zero-update free-rider, this time by more than $18$ points under the Dirichlet
partition and more than $5$ under the IID one. Thus, the main
patterns are consistent under selection-retrain and under online gating.
The last panel of Table~\ref{tab:online} adds selection-retrain under the IID
partition, whose Dirichlet counterpart appears in
Table~2 of the main paper.

\subsection{MNIST Repetition}
\label{app:mnist-rep}

Table~\ref{tab:mnist-int} repeats the same setting on MNIST. Outside the
gradient-noise column every row sits close to the ceiling, so the
methods differ less than on CIFAR-10. Nevertheless, the differentiating
conditions show the same pattern: the first-order methods remain at vanilla
performance under gradient noise, matching it exactly under selection-retrain
in the Dirichlet partition, and the renormalizing methods fall below vanilla
under the zero-update free-rider when the gate runs online. Under
selection-retrain the same methods return to vanilla's level in that condition
instead, so on MNIST the cost of mis-ranking a free-rider is visible only while
the gate is acting on the trajectory it is scoring.

\subsection{Exclusions of the LLM Gate}
\label{app:gate-sets}

Table~\ref{tab:gate-pr} evaluates the exclusions made by each gate under
selection-retrain and under online gating, rather than their downstream
performance. Precision is the fraction
of excluded units that are corrupted, recall is the fraction of corrupted
units excluded, and the false-exclusion rate is the fraction of clean units
excluded. A unit is a client in the left panel and a client--round pair after
the first 10 rounds in the right panel. The clean condition contains no corrupted clients,
so recall is undefined and precision is zero whenever the gate fires. We
therefore characterize that condition through the false-exclusion rate alone.
Under the two corruption conditions, the same quantity measures the cost to
clean clients. Because a gate can lower this rate simply by firing less, the
table ranks precision and recall but not the false-exclusion rate.

The left panel evaluates the gate whose EM is reported in
Table~3 of the main paper.
Neither variant excludes any client in the clean condition, so precision and
recall are undefined there as well. Under the zero-update free-rider, both
variants exclude exactly the 20 corrupted clients in every seed, which is why
their EM in Table~3 of the main paper matches oracle-exclusion. The two
corruption conditions
differ in kind rather than in degree. Under answer-swap, precision remains
perfect, so the gate never removes a clean client, but on average $5.0$ of the
20 corrupted clients survive for Flirds and $6.7$ for Flirds-1st. The
$1.5$-point gap to oracle-exclusion in
Table~3 of the main paper is
therefore entirely a recall gap, and the higher recall of Flirds is consistent
with its higher EM.

The right panel applies the same rule and the same gate parameters online, as
in Appendix~\ref{app:online}, and separates the two corruption
conditions much more sharply. A client counts as excluded in a round if the
gate held it out of the candidate pool or zeroed its weight. Once gating
begins, every client receives this status in every round, whether or not it was
sampled. Precision stays
high under both corruption conditions and the false-exclusion rate stays low,
so the exclusions the online gate makes are mostly correct. The difference
lies in coverage: it excludes the zero-update free-rider in nearly every
client-round, but reaches only a small fraction of
the answer-swap client-rounds, because the accumulated contribution of an
answer-swap client stays above the threshold for most of the trajectory. An
end-of-run gate can act on the evidence accumulated over the whole run,
whereas a per-round gate reading the running total largely cannot.

Online, the two variants also separate in a way they do not at the end of the
run: Flirds gates more aggressively and reaches more of the answer-swap
client-rounds. The right panel therefore trades coverage against precision
rather than ordering the two variants. Under the zero-update free-rider, their recall
instead coincides, so a free-rider is already visible in the first-order term,
which is consistent with the two variants sharing the zero-update column of
Table~3 of the main paper. It is the answer-swap client whose detection
depends on how the contribution is estimated.

\subsection{Removal Curve}
\label{app:removal}

Table~\ref{tab:removal-cnn} evaluates the contribution ranking using
retraining performance outside the attribution game \citep{jia2019towards}. A
method's ranking fixes two removal orders, one from the lowest rank upward and
one from the highest rank downward. At each removal point along an order we
retrain from the same initial weights with the clients that remain and record
test accuracy. This gives ten points per order at $N{=}10$. We compare the two
resulting curves point by point. A wider gap between them
indicates stronger alignment between the ranking and
performance. We report that gap averaged over all ten removal points. The
in-run Shapley and Flirds produce nearly identical gaps
under all three corruption conditions, differing by at most $0.15$~points.
Under gradient noise, only Flirds-1st reverses sign.
The zero-update free-rider divides the renormalizing methods instead of
sinking them uniformly: ShapleyFL reverses sign and FedSV is unstable across
seeds, whereas GTG-Shapley stays within $0.04$~points
of the in-run Shapley itself. This is narrower than the uniform failure of the
same methods under selection-retrain in
Section~5.3 of the main paper.
In the clean condition, every method is near zero, so removal
order has little effect on performance. This result is consistent with the
absence of a stable signal diagnosed in
Appendix~\ref{app:signal}.

% ===================================================================== F
\section{Computational Cost}
\label{app:cost}

Section~5.4 of the main paper reports contribution-evaluation cost in two
views: an analytic count of the dominant per-round operations and measured
wall-clock time. This appendix connects the two through a
per-operation ratio
and describes how the ratio between methods changes with the number of
participants per round.

\noindent\textbf{Per-operation microbenchmark.} Table~\ref{tab:microbench} times the two
operations that dominate the methods in Table~4 of the main paper: one
validation forward pass and one HVP along a prescribed direction. We use the same
software and hardware stack as the wall-clock measurements in that table. One
HVP costs $6.47$ validation forward passes. Both operations scale linearly
with the number of validation examples, so the ratio carries across validation
set sizes even though the individual times do not. This ratio compares the one
HVP per round used by Flirds with the $2^{|P_r|}$ forward passes per round used
by exhaustive enumeration, and yields the $6.5$ quoted in
Section~5.4 of the main paper.

\noindent\textbf{Reconciling the two views.} The wall-clock column evaluates every
method on one realized trajectory. Dividing the in-run Shapley total by its
operation count, $30 \times 2^{10}$ coalition evaluations, gives the cost of
one coalition evaluation in that run. This cost includes a validation forward
pass and coalition-model reconstruction. Multiplying it by the microbenchmark
ratio therefore gives an upper bound on one per-round HVP, and multiplying
again by 30 gives the predicted Flirds total in Table~\ref{tab:cost-check}. The
prediction exceeds the measurement by less than $0.5\%$, showing that the
operation counts and wall-clock times describe the same computation.

Two independent routes therefore reach the same ratio between Flirds and
exhaustive enumeration. From the operation counts, the in-run Shapley needs
$2^{10}$ forward passes per round against one HVP for Flirds, giving
$2^{10}/6.47 = 158.3$. From the measurements alone, the two totals of
Table~4 of the main paper stand in the ratio $24{,}975/157.3 = 158.8$. The $1/159$
of Section~5.4 of the main paper is this last number, and the two routes agree to
better than $1\%$, so the figure does not depend on which view is used. The
cost of the second-order curvature term is the ratio of an HVP to a validation
gradient, which the two Flirds rows of the same table measure directly as
$157.3/53.8 = 2.9$.

\noindent\textbf{Dependence on the number of participants.} Flirds uses one HVP per
round independently of $|P_r|$, plus $|P_r|$ inner products. Exhaustive
enumeration of the round game instead uses $2^{|P_r|}$ forward passes. Under
the dominant-operation approximation and the microbenchmark ratio, their cost
ratio grows as $2^{|P_r|}/6.47$: about $5$
at $|P_r|{=}5$, $158$ at $|P_r|{=}10$, and $5.1\times10^3$ at $|P_r|{=}15$.
Only the $|P_r|{=}10$ point is measured, and the other two are projections
from the operation counts and Table~\ref{tab:microbench} rather than separate
runs. GTG-Shapley, FedSV, and ComFedSV do not follow this curve, because their
permutation budgets are set by the convergence criteria of
Appendix~\ref{app:train} rather than by $|P_r|$ alone.

\begin{table}[t]
\centering
\small
\setlength{\tabcolsep}{4pt}
% [M1] per-op 마이크로벤치. 출처 rundir: runs/measured_2026-07/microbench/
% summary.json (forward_fp32 = 1.6012, hvp_fp32 = 10.3600). fp32, B200 1장,
% Llama-3.2-1B-Instruct, 10 chunks. 두 연산 모두 val 예제 수에 선형이라 비
% (6.47)만 val-불변이고, 절대 시간은 이 측정 조건에 묶인다 -- 다른 val 크기로
% 환산해 옮기지 말 것.
\begin{tabular}{lc}
\hline
Operation & Time (s) \\
\hline
validation forward pass & $1.601$ \\
one HVP & $10.360$ \\
\hline
HVP\,/\,forward pass & $6.47$ \\
\hline
\end{tabular}
\caption{Time per operation in seconds on the software and hardware stack used
for the wall-clock measurements in Table~4 of the main paper.}
\label{tab:microbench}
\end{table}

\begin{table}[t]
\centering
\small
\setlength{\tabcolsep}{4pt}
% [M2] op-count x 실측 교차검증. 예측 열 = 런 자체에서 얻은 coalition 평가
% 단가 24,975/(30 x 2^10) = 0.8130s 에 마이크로벤치 비 6.47 과 30 라운드를 곱한
% 값(157.8). 비 행은 2^10/6.47 vs 24,975/157.3. 실측 값은 본문 Table 4 그대로
% 인용 -- 여기서 자체 로그 값으로 바꾸면 본문과 어긋나므로 금지.
\begin{tabular}{lcc}
\hline
 & Operation counts & Measurement \\
\hline
Flirds total (s) & $157.8$ & $157.3$ \\
in-run SV\,/\,Flirds & $158.3$ & $158.8$ \\
\hline
\end{tabular}
\caption{Predicted and measured valuation cost over a complete
run (five domains, Dirichlet($\alpha{=}0.5$), Llama-3.2-1B-Instruct; $N{=}100$,
10/100 participation, $R{=}30$). The first row is a total in seconds. The
second row is the ratio of the two methods.}
\label{tab:cost-check}
\end{table}

% References (main.tex 와 같은 폴더의 references.bib 를 공유; \bibliographystyle 은
% aaai2027.sty 가 내부 설정. supplement 는 별도 컴파일이므로 자체 서지가 필요하다.
% 합본할 때는 위 병합 규약대로 이 \bibliography 와 \end{document} 를 옮기지 않는다.)
\bibliography{references}

\end{document}
